% L'exercice sur texte philosophique — Philosophie 1re Tech — Cours signé OnBuch+
% Programme officiel MINESEC. Compiler avec : tectonic N19-exercice-texte-philosophique.tex
\newcommand{\DOCMATIERE}{Philosophie}
\newcommand{\DOCNIVEAU}{1re Tech}
\newcommand{\DOCMODULE}{Philosophie – classe de Première technique}
\newcommand{\DOCLECON}{N19}
\newcommand{\DOCTITRE}{L'exercice sur texte philosophique}
% OnBuch+ — préambule « cours signés » ENRICHI (compilé avec tectonic / XeLaTeX)
% Système visuel « Pop » d'OnBuch+ : fond crème (jamais blanc), bordures encre
% épaisses, ombres « dures » décalées, titres Archivo Black en capitales.
% Version MATHÉMATIQUES : graphes (pgfplots/tikz), boîtes pédagogiques
% (définition, propriété, méthode, attention, le savais-tu, exercice résolu,
% exercice, TP, bilan), page de garde et signature OnBuch+ ancrée en bas.
%
% Macros attendues dans main.tex AVANT \input{preamble} :
%   \DOCMATIERE  \DOCNIVEAU  \DOCMODULE  \DOCLECON (numéro)  \DOCTITRE
\documentclass[11pt]{article}

\usepackage{fontspec}
\usepackage[french]{babel}
\usepackage[a4paper,margin=2cm,top=3.4cm,bottom=2.8cm,headheight=1.4cm,footskip=1.3cm]{geometry}
\usepackage{amsmath,amssymb,amsfonts}
\usepackage{mathtools} % \xmapsto, \coloneqq... souvent utilisés par les modèles
\usepackage{cancel} % simplifications barrées (bilans de forces)
% \abs{z} (module, valeur absolue) et \norm{u} (norme), utilisés par les modèles
\providecommand{\abs}{}\let\abs\relax\DeclarePairedDelimiter{\abs}{\lvert}{\rvert}
\providecommand{\norm}{}\let\norm\relax\DeclarePairedDelimiter{\norm}{\lVert}{\rVert}
\usepackage[table]{xcolor} % [table] : fournit aussi \rowcolors (alternance de lignes)
\usepackage{pagecolor}
\usepackage{tikz}
\usetikzlibrary{arrows.meta,positioning,calc,shapes.geometric,shapes.misc,shapes.arrows,decorations.pathmorphing,decorations.pathreplacing,decorations.markings,patterns,shadows,mindmap,trees,fit,backgrounds,matrix,intersections,angles,quotes}
\usepackage{pgfplots}
\usepackage[version=4]{mhchem} % équations nucléaires \ce{^{235}_{92}U}
\usepackage[european,siunitx]{circuitikz} % schémas électriques
\pgfplotsset{compat=1.18}
\usepgfplotslibrary{fillbetween,groupplots} % aires entre courbes (calcul intégral)
\usepackage{enumitem}
\usepackage{fancyhdr}
% [most] charge toutes les bibliothèques tcolorbox (listings, minted, raster,
% documentation...) et gonfle inutilement la mémoire TeX sur un document long
% (~300 boîtes) : on ne charge que ce qui est réellement utilisé.
\usepackage[skins,breakable]{tcolorbox}
\usepackage{graphicx}
\usepackage{colortbl}
\usepackage{booktabs}
\usepackage{tabularx}
\usepackage{multirow}
\usepackage{diagbox} % case d'en-tête diagonale des tableaux à double entrée
% Colonnes extensibles centrée / gauche / droite (C, L, R), souvent utilisées par les modèles
\newcolumntype{C}{>{\centering\arraybackslash}X}
\newcolumntype{L}{>{\raggedright\arraybackslash}X}
\newcolumntype{R}{>{\raggedleft\arraybackslash}X}
\usepackage{array}
\usepackage{multicol}
\usepackage{siunitx}
% parse-numbers=false : accepte aussi \SI{2,5 \times 10^{-3}}{...}
\sisetup{locale=FR,output-decimal-marker={,},group-minimum-digits=4,parse-numbers=false}
\DeclareSIUnit\ohm{\ensuremath{\symup{\Omega}}} % U+2126 absent de la police
\DeclareSIUnit\division{div} % réglages d'oscilloscope (ms/div, V/div)
\DeclareSIUnit\tr{tr} % tours (vitesse de rotation, ex. \SI{1500}{\tr\per\minute})
\DeclareSIUnit\molar{M} % concentration molaire (ex. \SI{3}{\molar}, \SI{1}{\micro\molar})
\DeclareSIUnit\an{an} % année (le modèle confond parfois avec \year, un compteur TeX) — ex. \SI{5}{\kilo\meter\per\an}
\DeclareSIUnit\FCFA{FCFA} % franc CFA (ex. \SI{500}{\FCFA\per\kilo\gram})
\DeclareSIUnit\degreeBrix{\textdegree Brix} % teneur en sucre d'un sirop/jus (réfractométrie)
\DeclareSIUnit\month{mois} % le modèle utilise parfois \month, un compteur TeX, comme unité de durée
\DeclareSIUnit\microliter{\micro\liter} % le modèle écrit parfois \microliter en un seul token
\DeclareSIUnit\million{millions} % dénombrement (ex. \SI{15}{\million\per\mL} spermatozoïdes)
\usepackage{qrcode}
% \degree : commande fréquemment utilisée par les modèles pour un angle ou une
% température en dehors de \SI{}{} (non fournie par les paquets chargés).
\providecommand{\degree}{^{\circ}}
% \qed (fin de démonstration), utilisé par les modèles sans amsthm
\providecommand{\qed}{\hfill$\square$}
% \barre : double barre des valeurs interdites dans un tableau de variations (tkz-tab)
\providecommand{\barre}{\ensuremath{\|}}
% environnement proof (amsthm non chargé) utilisé par les modèles
\ifdefined\proof\else\newenvironment{proof}[1][Démonstration]{\par\noindent\textit{#1.}\ }{\qed\par}\fi
\usepackage{lastpage}
\usepackage{hyperref}
\hypersetup{hidelinks}

% ============================================================
% Palette « Pop » OnBuch+ — mêmes hex que la classe P (plus_ui.dart)
% ============================================================
\definecolor{poporange}{HTML}{F59321}
\definecolor{popdark}{HTML}{E07A0C}
\definecolor{popcream}{HTML}{FFF6E8}
\definecolor{popink}{HTML}{1C1714}
\definecolor{poppurple}{HTML}{7A5AE0}
\definecolor{popgreen}{HTML}{2E9E6A}
\definecolor{poppink}{HTML}{E0457B}
\definecolor{popblue}{HTML}{2F7DE1}
\definecolor{popgold}{HTML}{C98A12}
\definecolor{popmuted}{HTML}{8A7F76}
\definecolor{popline}{HTML}{EADFD2}
\definecolor{popgray}{HTML}{8A7F76}
\colorlet{popgrey}{popgray}
% Alias tolérants (le modèle nomme parfois les couleurs différemment)
\colorlet{popwhite}{white}
\colorlet{popblack}{popink}
\colorlet{popred}{poppink}
\colorlet{popyellow}{popgold}
% Teintes claires (fonds de tableaux/encadrés) — le modèle nomme parfois
% les couleurs avec un suffixe L (ex. popblueL) sans les avoir définies.
\colorlet{poporangeL}{poporange!20}
\colorlet{popdarkL}{popdark!20}
\colorlet{poppurpleL}{poppurple!20}
\colorlet{popgreenL}{popgreen!20}
\colorlet{poppinkL}{poppink!20}
\colorlet{popblueL}{popblue!20}
\colorlet{popgoldL}{popgold!20}
\colorlet{popredL}{popred!20}
\colorlet{popgrayL}{popgray!20}
% Teintes claires pour fonds de tableaux / zones
\colorlet{popblueL}{popblue!10!white}
\colorlet{poporangeL}{poporange!14!white}
\colorlet{popgreenL}{popgreen!12!white}
\colorlet{poppurpleL}{poppurple!12!white}
\colorlet{poppinkL}{poppink!12!white}
\colorlet{popgoldL}{popgold!14!white}

\pagecolor{popcream}

% ============================================================
% Typographie
% ============================================================
\defaultfontfeatures{Ligatures=TeX}
\setmainfont{PlusJakartaSans-Medium.ttf}[Path=../../../fonts/,BoldFont=PlusJakartaSans-Bold.ttf]
% Maths en sans-serif (Fira Math) avec chiffres et lettres droites de Plus
% Jakarta, pour que les formules se fondent dans le texte.
\usepackage[math-style=ISO,bold-style=ISO]{unicode-math}
\setmathfont{FiraMath-Regular.otf}[Path=../../../fonts/]
\setmathfont{PlusJakartaSans-Medium.ttf}[Path=../../../fonts/,range={up/num,up/{Latin,latin},\mathup}]
\usepackage{icomma} % virgule décimale sans espace en mode math
% Caractères mathématiques Unicode absents des polices de texte (Plus Jakarta,
% Archivo Black) : rendus par leur équivalent mathématique, en texte comme en
% formule — sinon ils s'affichent en carré vide (ex. « Arithmétique dans ℤ »).
\usepackage{newunicodechar}
\newunicodechar{ℝ}{\ensuremath{\mathbb{R}}}
\newunicodechar{ℤ}{\ensuremath{\mathbb{Z}}}
\newunicodechar{ℂ}{\ensuremath{\mathbb{C}}}
\newunicodechar{ℕ}{\ensuremath{\mathbb{N}}}
\newunicodechar{ℚ}{\ensuremath{\mathbb{Q}}}
\newunicodechar{𝔻}{\ensuremath{\mathbb{D}}}
\newunicodechar{α}{\ensuremath{\alpha}}
\newunicodechar{β}{\ensuremath{\beta}}
\newunicodechar{γ}{\ensuremath{\gamma}}
\newunicodechar{δ}{\ensuremath{\delta}}
\newunicodechar{ε}{\ensuremath{\varepsilon}}
\newunicodechar{θ}{\ensuremath{\theta}}
\newunicodechar{λ}{\ensuremath{\lambda}}
\newunicodechar{μ}{\ensuremath{\mu}}
\newunicodechar{σ}{\ensuremath{\sigma}}
\newunicodechar{τ}{\ensuremath{\tau}}
\newunicodechar{φ}{\ensuremath{\varphi}}
\newunicodechar{ψ}{\ensuremath{\psi}}
\newunicodechar{ω}{\ensuremath{\omega}}
\newunicodechar{Σ}{\ensuremath{\Sigma}}
% majuscules produites par \MakeUppercase dans les titres (α-aminés -> Α)
\newunicodechar{Α}{\ensuremath{\alpha}}
\newunicodechar{Β}{\ensuremath{\beta}}
\newunicodechar{Γ}{\ensuremath{\Gamma}}
\newunicodechar{Δ}{\ensuremath{\Delta}}
\newunicodechar{Ε}{\ensuremath{\varepsilon}}
\newunicodechar{Θ}{\ensuremath{\Theta}}
\newunicodechar{Λ}{\ensuremath{\Lambda}}
\newunicodechar{Μ}{\ensuremath{\mu}}
\newunicodechar{Τ}{\ensuremath{\tau}}
\newunicodechar{Φ}{\ensuremath{\Phi}}
\newunicodechar{Ψ}{\ensuremath{\Psi}}
\newunicodechar{Ω}{\ensuremath{\Omega}}
\newunicodechar{↦}{\ensuremath{\mapsto}}
\newunicodechar{∈}{\ensuremath{\in}}
\newunicodechar{∉}{\ensuremath{\notin}}
\newunicodechar{∧}{\ensuremath{\wedge}}
\newunicodechar{⇔}{\ensuremath{\Leftrightarrow}}
\newunicodechar{⟺}{\ensuremath{\Longleftrightarrow}}
\newunicodechar{⇒}{\ensuremath{\Rightarrow}}
\newunicodechar{⟹}{\ensuremath{\Longrightarrow}}
\newunicodechar{∘}{\ensuremath{\circ}}
\newunicodechar{∪}{\ensuremath{\cup}}
\newunicodechar{∩}{\ensuremath{\cap}}
\newunicodechar{≡}{\ensuremath{\equiv}}
\newunicodechar{⊂}{\ensuremath{\subset}}
\newunicodechar{⊥}{\ensuremath{\perp}}
\newunicodechar{∃}{\ensuremath{\exists}}
\newunicodechar{∀}{\ensuremath{\forall}}
\newunicodechar{∛}{\ensuremath{\sqrt[3]{\,}}}
\newunicodechar{∜}{\ensuremath{\sqrt[4]{\,}}}
\newunicodechar{⃗}{\ensuremath{{}^{\to}}}
\newunicodechar{ⁿ}{\ifmmode^{n}\else\textsuperscript{\ensuremath{n}}\fi}
\newunicodechar{ˣ}{\ifmmode^{x}\else\textsuperscript{\ensuremath{x}}\fi}
\newunicodechar{ᵇ}{\ifmmode^{b}\else\textsuperscript{\ensuremath{b}}\fi}
\newunicodechar{ʳ}{\ifmmode^{r}\else\textsuperscript{\ensuremath{r}}\fi}
\newunicodechar{ᵘ}{\ifmmode^{u}\else\textsuperscript{\ensuremath{u}}\fi}
\newunicodechar{ʸ}{\ifmmode^{y}\else\textsuperscript{\ensuremath{y}}\fi}
\newunicodechar{ᵃ}{\ifmmode^{a}\else\textsuperscript{\ensuremath{a}}\fi}
\newunicodechar{ᵉ}{\ifmmode^{e}\else\textsuperscript{\ensuremath{e}}\fi}
\newunicodechar{ᵏ}{\ifmmode^{k}\else\textsuperscript{\ensuremath{k}}\fi}
\newunicodechar{ᵖ}{\ifmmode^{p}\else\textsuperscript{\ensuremath{p}}\fi}
\newunicodechar{ᴺ}{\ifmmode^{N}\else\textsuperscript{\ensuremath{N}}\fi}
\newunicodechar{ᶜ}{\ifmmode^{c}\else\textsuperscript{\ensuremath{c}}\fi}
\newunicodechar{ʰ}{\ifmmode^{h}\else\textsuperscript{\ensuremath{h}}\fi}
\newunicodechar{ᵠ}{\ifmmode^{q}\else\textsuperscript{\ensuremath{q}}\fi}
\newunicodechar{ₙ}{\ifmmode_{n}\else\textsubscript{\ensuremath{n}}\fi}
\newunicodechar{ₐ}{\ifmmode_{a}\else\textsubscript{\ensuremath{a}}\fi}
\newunicodechar{ᵢ}{\ifmmode_{i}\else\textsubscript{\ensuremath{i}}\fi}
\newunicodechar{ᵣ}{\ifmmode_{r}\else\textsubscript{\ensuremath{r}}\fi}
\newunicodechar{ₖ}{\ifmmode_{k}\else\textsubscript{\ensuremath{k}}\fi}
\newunicodechar{ᵦ}{\ifmmode_{\beta}\else\textsubscript{\ensuremath{\beta}}\fi}
\newunicodechar{ₓ}{\ifmmode_{x}\else\textsubscript{\ensuremath{x}}\fi}
\newfontfamily\popdisplay{ArchivoBlack-Regular.ttf}[Path=../../../fonts/]
\newfontfamily\poplogo{SpaceGrotesk-Bold.ttf}[Path=../../../fonts/]
\newfontfamily\popsemi{PlusJakartaSans-SemiBold.ttf}[Path=../../../fonts/]
\newfontfamily\popxbold{PlusJakartaSans-ExtraBold.ttf}[Path=../../../fonts/]
\newfontfamily\popxboldls{PlusJakartaSans-ExtraBold.ttf}[Path=../../../fonts/,LetterSpace=3.0]

\setlist[enumerate]{leftmargin=1.7em,itemsep=5pt,topsep=4pt}
\setlist[itemize]{leftmargin=1.7em,itemsep=4pt,topsep=4pt,label={\color{poporange}\textbullet}}
\setlength{\parindent}{0pt}
\setlength{\parskip}{5pt}
\renewcommand{\arraystretch}{1.25}

% pgfplots : style Pop par défaut
\pgfplotsset{
  every axis/.append style={axis line style={popink,line width=1pt},tick style={popink},
    grid style={popline,line width=0.5pt},label style={font=\small},tick label style={font=\footnotesize}},
  popcurve/.style={poporange,line width=1.8pt,no markers,smooth},
}
% Même style disponible dans un \draw TikZ simple (hors pgfplots)
\tikzset{popcurve/.style={poporange,line width=1.8pt}}
\tikzset{popgrid/.style={popline,very thin}}
\colorlet{popcurve}{poporange} % le nom du style est parfois utilisé comme couleur

% ============================================================
% Pill / squiggle
% ============================================================
\newcommand{\poppill}[3]{%
  \tikz[baseline=-0.55ex]\node[draw=popink,line width=1.2pt,rounded corners=2.6pt,%
    fill=#2,inner xsep=6.5pt,inner ysep=3pt]{%
    \popxboldls\fontsize{7}{7}\selectfont\color{#3}\MakeUppercase{#1}};%
}
\newcommand{\popsquiggle}[1]{%
  \tikz[baseline=-0.4ex]{
    \draw[#1,line width=2.6pt,line cap=round]
      plot [smooth] coordinates {(0,0.09) (0.34,0.01) (0.68,0.21) (1.02,0.01) (1.36,0.10)};
  }%
}
% Mot-clé surligné (marqueur orange sous le texte)
\newcommand{\cle}[1]{{\popxbold\color{popdark}#1}}

% ============================================================
% En-tête / pied
% ============================================================
\pagestyle{fancy}
\fancyhf{}
\renewcommand{\headrulewidth}{0pt}
\renewcommand{\footrulewidth}{0pt}
\lhead{\raisebox{-0.15ex}{\poplogo\fontsize{13}{13}\selectfont\color{popink}OnBuch\color{poporange}+}}
\rhead{\poppill{\DOCMATIERE\ \textbullet\ \DOCNIVEAU\ \textbullet\ Leçon \DOCLECON}{white}{popink}}
\lfoot{\popxbold\fontsize{7.6}{7.6}\selectfont\color{popmuted}\MakeUppercase{Cours signé OnBuch+}\ \textbullet\ {\popsemi\DOCTITRE}}
\rfoot{\tikz[baseline=-0.6ex]\node[rounded corners=8.5pt,fill=popink,minimum height=17pt,minimum width=17pt,inner xsep=5pt,inner sep=0]{\popxbold\fontsize{8}{8}\selectfont\color{popcream}\thepage\,/\,\pageref*{LastPage}};}

% ============================================================
% Titres
% ============================================================
\newcommand{\chaptitre}[2]{%
  \begin{tcolorbox}[enhanced,colback=poporange,colframe=popink,boxrule=2.6pt,arc=11pt,
      drop shadow=popink,left=15pt,right=15pt,top=12pt,bottom=11pt,before skip=3pt,after skip=18pt]
    \poppill{#2}{popink}{popcream}\par\vspace{9pt}
    {\raggedright\hyphenpenalty=10000\exhyphenpenalty=10000\popdisplay\fontsize{22}{24}\selectfont\color{popcream}\MakeUppercase{#1}\par}
  \end{tcolorbox}
}
% Section numérotée : pastille orange avec le numéro + titre Archivo
\newcounter{popsec}
\newcommand{\coursec}[1]{%
  \stepcounter{popsec}\setcounter{popsub}{0}%
  \par\vspace{16pt}\noindent
  \begin{minipage}{\linewidth}
  \tikz[baseline=-0.7ex]\node[draw=popink,line width=1.6pt,fill=poporange,rounded corners=4pt,minimum size=22pt,inner sep=0]{\popdisplay\fontsize{12}{12}\selectfont\color{popcream}\thepopsec};%
  \hspace{8pt}{\popdisplay\fontsize{15}{17}\selectfont\color{popink}\MakeUppercase{#1}}\par
  \vspace{4pt}\noindent{\color{poporange}\rule{36pt}{3.6pt}}
  \end{minipage}\par\nopagebreak\vspace{8pt}
}
\newcounter{popsub}
\newcommand{\courssub}[1]{%
  \stepcounter{popsub}%
  \par\vspace{9pt}\noindent{\popxbold\fontsize{11.5}{13}\selectfont\color{popink}{\color{poporange}\thepopsec.\thepopsub}\hspace{6pt}#1}\par\nopagebreak\vspace{4pt}
}

% ============================================================
% Boîtes pédagogiques — même recette : fond blanc, bordure épaisse colorée,
% ombre dure encre, badge plein. Titre optionnel [#1].
% ============================================================
\tcbset{popbase/.style={breakable,enhanced,colback=white,boxrule=2.3pt,arc=9pt,
  shadow={3pt}{-3pt}{0pt}{popink},left=14pt,right=14pt,top=11pt,bottom=11pt,before skip=11pt,after skip=15pt,
  fontupper=\popsemi\fontsize{10.6}{14.5}\selectfont\color{popink},
  fontlower=\popsemi\fontsize{10.6}{14.5}\selectfont\color{popink}}}
% Coche dessinée (le glyphe ✓ n'existe pas dans les polices du document)
\newcommand{\popcheck}[1]{\tikz[baseline=-0.5ex]\draw[#1,line width=1.6pt,line cap=round,line join=round](-0.13,0.0)--(-0.04,-0.09)--(0.13,0.1);}
\providecommand{\coche}{\popcheck{popgreen}}
\providecommand{\resultat}[1]{\par\smallskip\noindent\fcolorbox{popgreen}{popgreenL}{\parbox{\dimexpr\linewidth-2\fboxsep-2\fboxrule\relax}{\textbf{Résultat.} #1}}\par\smallskip}
\newcommand{\popboxhead}[3]{\poppill{#1}{#2}{white}\ifx\relax#3\relax\else\hspace{6pt}{\popxbold\fontsize{10.6}{12}\selectfont\color{popink}#3}\fi\par\vspace{7pt}}

\newtcolorbox{aretenir}[1][]{popbase,colframe=popgreen,before upper={\popboxhead{À retenir}{popgreen}{#1}}}
\newtcolorbox{exemplebox}[1][]{popbase,colframe=poporange,before upper={\popboxhead{Exemple}{poporange}{#1}}}
\newtcolorbox{definition}[1][]{popbase,colframe=popblue,before upper={\popboxhead{Définition}{popblue}{#1}}}
\newtcolorbox{propriete}[1][]{popbase,colframe=poppurple,before upper={\popboxhead{Propriété}{poppurple}{#1}}}
\newtcolorbox{methode}[1][]{popbase,colframe=popgold,before upper={\popboxhead{Méthode}{popgold}{#1}}}
\newtcolorbox{attention}[1][]{popbase,colframe=poppink,before upper={\popboxhead{Attention}{poppink}{#1}}}
\newtcolorbox{experience}[1][]{popbase,colframe=popblue,colback=popblueL,before upper={\popboxhead{Expérience}{popblue}{#1}}}
% « Le savais-tu ? » : carte violette pleine (contexte, culture, Cameroun)
\newtcolorbox{savaistu}[1][]{popbase,colback=poppurple,colframe=popink,
  fontupper=\popsemi\fontsize{10.4}{14.2}\selectfont\color{white},
  before upper={\poppill{Le savais-tu\,?}{white}{poppurple}\ifx\relax#1\relax\else\hspace{6pt}{\popxbold\color{white}#1}\fi\par\vspace{7pt}}}
% Exercice résolu : énoncé en haut, \tcblower puis la solution sur fond crème
\newcounter{popexo}\newcounter{popexr}
\newtcolorbox{exoresolu}[1][]{popbase,colframe=poporange,colbacklower=popcream,
  segmentation style={popink,line width=1.2pt,dashed},
  before upper={\stepcounter{popexr}\popboxhead{Exercice résolu \thepopexr}{poporange}{#1}},
  before lower={\poppill{Solution}{popink}{popcream}\par\vspace{6pt}}}
% Exercice (non résolu en place) — la correction est dans la partie Corrigés
\newtcolorbox{exercice}[1][]{popbase,colframe=popink,
  before upper={\stepcounter{popexo}\popboxhead{Exercice \thepopexo}{popink}{#1}}}
% Corrigé
\newtcolorbox{corrige}[1][]{popbase,colframe=popgreen,colback=popgreenL,
  before upper={\popboxhead{Corrigé}{popgreen}{#1}}}
% Objectifs / prérequis / bilan
\newtcolorbox{objectifs}{popbase,colframe=popink,colback=poporangeL,
  before upper={\popboxhead{Objectifs de la leçon}{popink}{}}}
\newtcolorbox{prerequis}{popbase,colframe=popink,colback=popblueL,
  before upper={\popboxhead{Prérequis}{popblue}{}}}
\newtcolorbox{bilan}{popbase,colframe=popink,colback=white,boxrule=2.8pt,
  before upper={\popboxhead{Fiche bilan}{poporange}{L'essentiel à connaître pour l'examen}}}
% Figure encadrée avec légende
\newtcolorbox{popfigure}[1][]{enhanced,breakable=false,colback=white,colframe=popline,boxrule=1.4pt,arc=8pt,
  left=10pt,right=10pt,top=10pt,bottom=8pt,before skip=10pt,after skip=12pt,halign=center,
  lower separated=false,halign lower=center,
  fontlower=\popsemi\fontsize{9}{11.5}\selectfont\color{popmuted},
  #1}
\newcommand{\legende}[1]{\par\vspace{6pt}{\popsemi\fontsize{9}{11.5}\selectfont\color{popmuted}#1}}

% ============================================================
% Signature OnBuch+
% ============================================================
\newcommand{\poplogoinline}[1]{{\poplogo\fontsize{#1}{#1}\selectfont\color{popink}OnBuch\color{poporange}+}}
\newcommand{\signatureonbuch}{%
  \begin{tcolorbox}[enhanced,colback=popink,colframe=popink,boxrule=2.2pt,arc=10pt,
      drop shadow=poporange,left=14pt,right=14pt,top=10pt,bottom=10pt,before skip=0pt,after skip=0pt]
    \tikz[baseline=-0.7ex]\node[circle,fill=popgreen,draw=popcream,line width=1.3pt,minimum size=22pt,inner sep=0]{\popcheck{white}};%
    \hspace{9pt}\begin{minipage}[c]{0.62\linewidth}
      {\popxbold\fontsize{12}{13}\selectfont\color{popcream}Signé\ \textbullet\ {\poplogo OnBuch\color{poporange}+}}\par\vspace{2pt}
      {\raggedright\popsemi\fontsize{8}{10.5}\selectfont\color{popline}\MakeUppercase{Équipe pédagogique OnBuch+}\\\MakeUppercase{Conforme au programme officiel MINESEC}\par}
    \end{minipage}\hfill
    {\poplogo\fontsize{22}{22}\selectfont\color{popcream}OnBuch\color{poporange}+}
  \end{tcolorbox}%
}

% ============================================================
% Page de garde
% #1 titre · #2 matière · #3 niveau · #4 module · #5 n° leçon · #6 durée ·
% #7 description / objectifs (texte ou itemize)
% ============================================================
\newcommand{\pagegarde}[7]{%
  \thispagestyle{empty}
  \newgeometry{a4paper,margin=1.7cm,top=1.6cm,bottom=1.4cm}%
  \begin{tikzpicture}[remember picture,overlay]
    \fill[poporange!18] ([xshift=-3.2cm,yshift=-3.6cm]current page.north east) circle (4.6cm);
    \fill[poppurple!14] ([xshift=2.4cm,yshift=7.6cm]current page.south west) circle (3.8cm);
  \end{tikzpicture}%
  {\poplogo\fontsize{30}{30}\selectfont\color{popink}OnBuch\color{poporange}+}\hfill
  \raisebox{0.6ex}{\poppill{Cours signé \textbullet\ Premium}{popink}{popcream}}\par
  \vspace{1pt}\popsquiggle{poppurple}\par
  \vspace{8pt}
  \begin{tcolorbox}[enhanced,colback=poporange,colframe=popink,boxrule=2.8pt,arc=13pt,
      drop shadow=popink,left=18pt,right=18pt,top=12pt,bottom=13pt,after skip=10pt]
    \poppill{#2 \textbullet\ #3}{popink}{popcream}\hspace{5pt}\poppill{#4}{white}{popink}\par\vspace{8pt}
    {\popdisplay\fontsize{14}{14}\selectfont\color{popink}LEÇON #5}\par\vspace{4pt}
    {\raggedright\hyphenpenalty=10000\exhyphenpenalty=10000\popdisplay\fontsize{25}{27}\selectfont\color{popcream}\MakeUppercase{#1}\par}
    \vspace{8pt}
    {\popxbold\fontsize{9.5}{9.5}\selectfont\color{popink}\MakeUppercase{Durée conseillée : #6}}
  \end{tcolorbox}
  \begin{tcolorbox}[enhanced,colback=white,colframe=popink,boxrule=2.2pt,arc=10pt,
      drop shadow=popink,left=16pt,right=16pt,top=10pt,bottom=10pt,after skip=8pt]
    \poppill{Dans cette leçon}{poppurple}{white}\par\vspace{5pt}
    \popsemi\fontsize{9.3}{12.6}\selectfont\color{popink}#7
  \end{tcolorbox}
  \noindent\begin{minipage}[t]{0.487\linewidth}
    \begin{tcolorbox}[enhanced,colback=popgreen,colframe=popink,boxrule=2.2pt,arc=10pt,
        drop shadow=popink,left=11pt,right=11pt,top=10pt,bottom=10pt,height=3.1cm,valign=center]
      \tcbox[colback=white,colframe=popink,boxrule=1.3pt,arc=5pt,boxsep=0pt,nobeforeafter,
        left=3pt,right=3pt,top=3pt,bottom=3pt,valign=center]{\qrcode[height=1.9cm]{https://play.google.com/store/apps/details?id=com.luvvix.onbuch&hl=fr}}%
      \hspace{8pt}\begin{minipage}[c]{0.5\linewidth}
        {\popdisplay\fontsize{11}{12.5}\selectfont\color{white}\MakeUppercase{Télécharge OnBuch}}\par\vspace{4pt}
        {\raggedright\popsemi\fontsize{8.4}{11}\selectfont\color{white}Gratuit \textbullet\ Google Play\\Tuteur IA, annales, concours, résultats\par}
      \end{minipage}
    \end{tcolorbox}
  \end{minipage}\hfill
  \begin{minipage}[t]{0.487\linewidth}
    \begin{tcolorbox}[enhanced,colback=poppurple,colframe=popink,boxrule=2.2pt,arc=10pt,
        drop shadow=popink,left=11pt,right=11pt,top=10pt,bottom=10pt,height=3.1cm,valign=center]
      \tikz[baseline=-0.7ex]\node[circle,fill=white,draw=popink,line width=1.3pt,minimum size=1.6cm,inner sep=0]{\popdisplay\fontsize{24}{24}\selectfont\color{poppurple}+};%
      \hspace{8pt}\begin{minipage}[c]{0.6\linewidth}
        {\popdisplay\fontsize{11}{12.5}\selectfont\color{white}\MakeUppercase{Passe à OnBuch+}}\par\vspace{4pt}
        {\raggedright\popsemi\fontsize{8.4}{11}\selectfont\color{white}Tous les cours signés, Tuteur IA illimité, fiches PDF — onglet OnBuch+ de l'app\par}
      \end{minipage}
    \end{tcolorbox}
  \end{minipage}
  \vspace{8pt}
  \popcontenu
  \vfill
  \signatureonbuch
  \restoregeometry
  \newpage
}

% Rangée « Ce cours contient » (page de garde)
\newcommand{\poptile}[3]{% #1 couleur #2 gros texte #3 légende
  \begin{tcolorbox}[enhanced,colback=#1,colframe=popink,boxrule=2pt,arc=9pt,shadow={2.5pt}{-2.5pt}{0pt}{popink},
      left=6pt,right=6pt,top=9pt,bottom=9pt,halign=center,valign=center,height=1.75cm,width=\linewidth,nobeforeafter]
    {\popdisplay\fontsize{12.5}{13}\selectfont\color{white}\MakeUppercase{#2}}\par\vspace{3pt}
    {\centering\popsemi\fontsize{7.8}{9.5}\selectfont\color{white}#3\par}
  \end{tcolorbox}}
\newcommand{\popcontenu}{%
  \par\noindent{\popxboldls\fontsize{7.5}{7.5}\selectfont\color{popmuted}CE COURS SIGNÉ CONTIENT}\par\vspace{7pt}
  \noindent\begin{minipage}[t]{0.235\linewidth}\poptile{popblue}{Cours}{complet, illustré, pas à pas}\end{minipage}\hfill
  \begin{minipage}[t]{0.235\linewidth}\poptile{poporange}{Méthodes}{et exercices résolus}\end{minipage}\hfill
  \begin{minipage}[t]{0.235\linewidth}\poptile{poppink}{Exercices}{type examen + corrigés}\end{minipage}\hfill
  \begin{minipage}[t]{0.235\linewidth}\poptile{popgreen}{Fiche bilan}{l'essentiel à retenir}\end{minipage}\par}

% Sommaire
\newtcolorbox{sommaire}{enhanced,colback=white,colframe=popink,boxrule=2.2pt,arc=10pt,
  drop shadow=popink,left=18pt,right=18pt,top=13pt,bottom=13pt,before skip=6pt,after skip=20pt,
  before upper={\poppill{Sommaire}{poppurple}{white}\par\vspace{9pt}\popsemi\fontsize{10.6}{14.5}\selectfont\color{popink}}}

% Signature de fin de cours — poussée tout en bas de la dernière page
\newcommand{\findecours}{%
  \par\vfill\nopagebreak
  \begin{center}\popsquiggle{poporange}\end{center}\vspace{4pt}
  \signatureonbuch
}
\AtBeginDocument{\def\llbracket{\mathopen{[\mkern-3.5mu[}}\def\rrbracket{\mathclose{]\mkern-3.5mu]}}} % intervalles d'entiers (glyphe ⟦ absent de la police)
% Glyphes absents de Fira Math : redéfinis à partir de glyphes disponibles
\AtBeginDocument{%
  \def\setminus{\mathbin{\backslash}}\let\smallsetminus\setminus
  \def\vdots{\mathord{\vbox{\baselineskip3.2pt\lineskiplimit0pt\kern4pt\hbox{.}\hbox{.}\hbox{.}}}}%
  \def\ddots{\mathinner{\mkern1mu\raise6.5pt\hbox{.}\mkern2mu\raise3.6pt\hbox{.}\mkern2mu\raise0.7pt\hbox{.}\mkern1mu}}%
  \def\triangle{\mathord{\scalebox{0.9}{$\symup{\Delta}$}}}%
  \def\nabla{\mathord{\raisebox{\depth}{\scalebox{1}[-1]{$\symup{\Delta}$}}}}%
  \def\checkmark{\popcheck{popdark}}%
  \def\star{\mathbin{\ast}}\def\bigstar{\mathord{\ast}}%
  \def\diamond{\mathbin{\scalebox{0.6}{$\Diamond$}}}%
  \def\bigcup{\mathop{\mathchoice{\raisebox{-0.3em}{\scalebox{1.7}{$\cup$}}}{\scalebox{1.25}{$\cup$}}{\cup}{\cup}}\displaylimits}%
  \def\bigcap{\mathop{\mathchoice{\raisebox{-0.3em}{\scalebox{1.7}{$\cap$}}}{\scalebox{1.25}{$\cap$}}{\cap}{\cap}}\displaylimits}%
  \def\lesseqqgtr{\mathrel{\lessgtr}}\def\bigoplus{\mathop{\oplus}\displaylimits}%
}
\providecommand{\ssi}{si et seulement si} % abréviation fréquente
\AtBeginDocument{\providecommand{\wideparen}{\overparen}} % arc de cercle

%% FIGFIX-BEGIN v4 — correctifs globaux des figures (docs/onbuch-generation/tools/figfix)
\usepackage{adjustbox}\usepackage{etoolbox}
% toute figure TikZ/pgfplots plus large que la ligne est réduite à la largeur disponible
\BeforeBeginEnvironment{tikzpicture}{\begin{adjustbox}{max width=\linewidth}}
\AfterEndEnvironment{tikzpicture}{\end{adjustbox}}
% légendes pgfplots lisibles par-dessus les courbes
\pgfplotsset{every axis/.append style={legend style={fill=white,fill opacity=0.92,text opacity=1,draw=black!15,inner sep=3pt}}}
% étiquettes d'arêtes (syntaxe edge["..."]) avec fond blanc
\tikzset{every edge quotes/.append style={fill=white,inner sep=1.2pt}}
%% FIGFIX-END
\begin{document}

\pagegarde{L'exercice sur texte philosophique}{Philosophie}{1re Tech}{Philosophie – classe de Première technique}{N19}{3 séances (fiche de progression harmonisée)}{Cette leçon outille l'élève de Première Technique pour maîtriser l'exercice canonique du baccalauréat : l'explication de texte philosophique, de l'analyse des concepts à la rédaction de l'essai personnel, en passant par les réponses aux questions guidées.
\vspace{4pt}
\begin{multicols}{2}\small
\begin{itemize}
\item À la fin de cette leçon, je sais identifier la structure logique (thèse, arguments, concepts) d'un texte philosophique quelconque.
\item À la fin de cette leçon, je sais distinguer les trois temps de l'exercice : explication, réponses aux questions, essai personnel.
\item À la fin de cette leçon, je sais appliquer la méthode des "trois lectures" pour dégager la problématique implicite d'un auteur.
\item À la fin de cette leçon, je sais répondre aux questions de compréhension et d'analyse en mobilisant mes connaissances du programme sans hors-sujet.
\item À la fin de cette leçon, je sais construire un plan dialectique ou analytique pour l'essai personnel en lien avec le texte.
\item À la fin de cette leçon, je sais rédiger une introduction (accroche, sujet, thèse, plan) et une conclusion (bilan, ouverture) conformes aux attendus du Bac.
\end{itemize}\end{multicols}}

\begin{sommaire}
\begin{multicols}{2}
\begin{enumerate}
\item 1. Généralités : Nature et enjeux de l'exercice
\item 2. Méthodologie de l'explication de texte (Étape A)
\item 3. Réponses aux questions guidées (Étape B)
\item 4. L'essai personnel : De la problématisation à la rédaction (Étape C)
\item 5. TP Intégré : Simulation Bac sur corpus 'Philosophie et Développement au Cameroun'
\item 6. Erreurs fréquentes, Remédiation et Excellence
\item Activité d'intégration
\item Méthodes et exercices résolus
\item Exercices
\item Corrigés des exercices
\item Fiche bilan
\end{enumerate}
\end{multicols}
\end{sommaire}

\begin{objectifs}
\begin{itemize}
\item À la fin de cette leçon, je sais identifier la structure logique (thèse, arguments, concepts) d'un texte philosophique quelconque.
\item À la fin de cette leçon, je sais distinguer les trois temps de l'exercice : explication, réponses aux questions, essai personnel.
\item À la fin de cette leçon, je sais appliquer la méthode des "trois lectures" pour dégager la problématique implicite d'un auteur.
\item À la fin de cette leçon, je sais répondre aux questions de compréhension et d'analyse en mobilisant mes connaissances du programme sans hors-sujet.
\item À la fin de cette leçon, je sais construire un plan dialectique ou analytique pour l'essai personnel en lien avec le texte.
\item À la fin de cette leçon, je sais rédiger une introduction (accroche, sujet, thèse, plan) et une conclusion (bilan, ouverture) conformes aux attendus du Bac.
\item À la fin de cette leçon, je sais utiliser les connecteurs logiques pour assurer la cohérence de mon argumentation écrite.
\item À la fin de cette leçon, je sais auto-évaluer ma copie grâce à la grille de correction type MINESEC (respect du texte, rigueur, culture).
\end{itemize}
\end{objectifs}

\begin{prerequis}
\begin{itemize}
\item Maîtrise de la lecture compréhensive (niveau 4ème/3ème) : repérer l'idée principale d'un paragraphe.
\item Connaissance des principaux concepts du programme de Terminale (Liberté, Conscience, Politique, Savoir, Vérité).
\item Utilisation basique des connecteurs logiques (donc, cependant, par conséquent, en outre).
\item Structure de base de la dissertation philosophique (Introduction, Développement en 2-3 parties, Conclusion).
\item Culture générale minimale sur quelques auteurs au programme (Descartes, Kant, Marx, Hountondji, Wiredu, Mbembe).
\end{itemize}
\end{prerequis}

\newpage

\coursec{Généralités : Nature et enjeux de l'exercice}

\textbf{Situation d'accroche.} Au lycée de Nkolbisson, Armande fixe le sujet du Bac blanc : « Expliquez le texte suivant de Paulin Hountondji sur la philosophie africaine. » Elle sait lire. Elle a même déjà entendu parler de Hountondji en cours. Mais devant la feuille blanche, une angoisse monte : comment transformer sa compréhension, encore floue et intuitive, en une copie structurée, rigoureuse, qui rapporte des points ? Faut-il donner son avis ? Faut-il résumer ? Faut-il critiquer l'auteur ? Combien de pages ? Combien de temps ?

Cette situation est celle de milliers de candidats chaque année au Cameroun. La bonne nouvelle, c'est que l'exercice sur texte philosophique n'est pas une épreuve d'inspiration : c'est une épreuve de \cle{méthode}. Comme un technicien qui suit une procédure de montage, le candidat qui connaît les étapes exactes peut dompter n'importe quel texte, du sujet d'examen aux débats de la vie courante (un discours politique, un éditorial de journal, un sermon). Cette leçon donne cette méthode.

\textbf{Problématique.} Qu'est-ce qu'un texte philosophique, en quoi l'exercice sur texte diffère-t-il de la dissertation, et comment est-il structuré et noté au Probatoire et au Baccalauréat technique ?

\courssub{Qu'est-ce qu'un texte philosophique ?}

Commençons par l'intuition. Tous les textes ne se valent pas. Un mode d'emploi de téléphone \emph{informe}, un poème \emph{émeut}, un récit \emph{raconte}. Un texte philosophique, lui, fait quelque chose de très précis : il \emph{démontre}. Son auteur ne se contente pas d'affirmer une opinion ; il la justifie par des raisonnements que le lecteur peut suivre, vérifier et discuter.

\begin{definition}[Texte philosophique]
Un \cle{texte philosophique} est un écrit argumenté qui vise à démontrer une \cle{thèse} (une position soutenue sur un problème) par la seule force de la raison, en utilisant des \cle{concepts} précis (des idées générales clairement définies, comme « liberté », « devoir », « vérité », « développement »).
\end{definition}

Trois marques permettent de reconnaître un texte philosophique :

\begin{itemize}
\item \textbf{Un problème} : le texte part d'une question ou d'une difficulté (ex. : la philosophie africaine existe-t-elle ? La tradition est-elle un obstacle au développement ?).
\item \textbf{Une thèse} : l'auteur prend position et défend une réponse déterminée.
\item \textbf{Des arguments} : l'auteur enchaîne des raisons, des distinctions conceptuelles, des exemples, pour convaincre rationnellement — et non par l'émotion, l'autorité ou la menace.
\end{itemize}

\begin{exemplebox}[Observer pour comprendre : démarche d'investigation]
Prenons deux phrases et comparons-les comme un technicien compare deux matériaux.
\begin{itemize}
\item Phrase A : « La tradition, c'est sacré, il faut la respecter ! »
\item Phrase B : « Si par tradition on entend la transmission critique d'un héritage, alors elle n'exclut pas l'innovation ; elle en est au contraire la condition, car innover suppose de partir d'un acquis. »
\end{itemize}
Observation : la phrase A \emph{affirme} et fait appel au sentiment ; la phrase B \emph{définit} un concept (« tradition »), pose une \emph{hypothèse} (« si... alors ») et tire une \emph{conséquence}. Seule la phrase B est philosophique. Conclusion de l'investigation : le critère du texte philosophique n'est pas le sujet traité, mais la \emph{manière rationnelle} de le traiter.
\end{exemplebox}

Au programme de Première technique, les textes proposés à l'examen sont des extraits courts (10 à 25 lignes environ) d'œuvres classiques ou contemporaines. Au Cameroun, le corpus inclut fréquemment des auteurs de la philosophie africaine, en lien avec les unités du programme (philosophie et développement, culture, technique, vérité).

\begin{savaistu}[Le corpus camerounais]
Depuis la réforme de 2018, les sujets du Baccalauréat technique au Cameroun intègrent régulièrement des auteurs de la philosophie africaine : Paulin Hountondji (critique de l'ethnophilosophie), Valentin-Yves Mudimbe (l'invention de l'Afrique), ou des écrivains penseurs comme Guillaume Oyono-Mbia, dont les pièces (\emph{Trois prétendants... un mari}) posent de vrais problèmes philosophiques sur la tradition et la modernité. Lire ces auteurs pendant l'année, c'est préparer l'examen \emph{et} nourrir sa culture personnelle.
\end{savaistu}

\courssub{Une épreuve obligatoire au Probatoire et au Baccalauréat technique}

\begin{definition}[L'exercice sur texte au programme MINESEC]
L'\cle{exercice sur texte philosophique} (ou explication de texte) est une épreuve obligatoire de l'examen philosophique (Probatoire en fin de Première technique, Baccalauréat technique en Terminale), au coefficient variable selon les séries. Le candidat reçoit un extrait d'œuvre philosophique et doit produire une copie en \textbf{trois parties distinctes} : l'explication du texte, les réponses aux questions posées, et un essai personnel.
\end{definition}

Cette épreuve évalue cinq compétences précises, que le correcteur recherche explicitement dans la copie :

\begin{enumerate}
\item \textbf{La compréhension} : avez-vous saisi le sens global du texte, son problème et sa thèse ?
\item \textbf{L'analyse conceptuelle} : savez-vous dégager et définir les concepts clés, suivre l'articulation des arguments ?
\item \textbf{L'argumentation} : savez-vous raisonner à votre tour, de façon ordonnée et justifiée ?
\item \textbf{La culture philosophique} : savez-vous mobiliser les notions de l'unité et les auteurs étudiés en classe ?
\item \textbf{L'expression française} : clarté, correction grammaticale et orthographique, vocabulaire philosophique approprié.
\end{enumerate}

\begin{exemplebox}[Un exemple chiffré : les statistiques parlent]
Données indicatives du Baccalauréat technique 2023 : moyenne nationale en philosophie de 9,2/20, avec un écart type de 3,5 points. Cet écart type élevé signifie que les copies sont très dispersées : l'épreuve \emph{discrimine} fortement. L'analyse des copies montre un fait décisif : les copies dépassant 14/20 sont celles qui maîtrisent la structure en trois parties distinctes (explication, questions, essai), clairement annoncées et séparées. Autrement dit, la méthode seule fait gagner plusieurs points, indépendamment du « génie » du candidat.
\end{exemplebox}

\courssub{Explication de texte et dissertation : deux exercices fondamentalement différents}

C'est la distinction la plus importante de cette leçon. Beaucoup d'élèves échouent parce qu'ils confondent les deux exercices.

\begin{itemize}
\item Dans la \cle{dissertation}, le sujet est une \emph{question} (« La technique libère-t-elle l'homme ? »). C'est \textbf{vous} qui construisez la réponse : vous problématisez, vous confrontez des thèses, vous concluez. L'exercice est une \emph{création personnelle} de pensée.
\item Dans l'\cle{explication de texte}, la réponse est \emph{déjà là}, dans le texte. Votre travail n'est pas de dire ce que \emph{vous} pensez du thème, mais de montrer ce que \emph{l'auteur} pense, comment il le prouve, et pourquoi son raisonnement est construit ainsi. L'exercice exige la \emph{fidélité à l'auteur}.
\end{itemize}

\begin{attention}[Le piège mortel : le hors-sujet]
L'erreur la plus fréquente et la plus sanctionnée : faire une dissertation \emph{sur le thème} du texte au lieu d'\emph{expliquer le texte} lui-même. Si le texte de Hountondji critique l'ethnophilosophie, le candidat qui écrit trois pages générales « sur la philosophie africaine » sans suivre l'argumentation précise de l'auteur fait du hors-sujet. Le correcteur le sanctionne lourdement, même si la copie est brillante. Règle d'or : dans la partie A, chaque affirmation doit pouvoir être rattachée à une ligne précise du texte. Le « je pense que... » n'a sa place que dans l'essai personnel (partie C).
\end{attention}

\begin{popfigure}
\begin{center}
\begin{tabularx}{\linewidth}{|l|X|X|}
\hline
\rowcolor{popblueL} \textbf{Critère} & \textbf{Explication de texte} & \textbf{Dissertation} \\
\hline
\textbf{Objectif} & Restituer et éclaircir la pensée de l'auteur : problème, thèse, arguments, enjeux. & Construire sa propre réponse à une question problématique. \\
\hline
\textbf{Posture} & Fidélité : on se met « au service » du texte, on le suit pas à pas. & Liberté : on dialogue avec les auteurs, on confronte les thèses. \\
\hline
\textbf{Structure} & Suit l'ordre et les étapes du texte (linéaire) ou sa logique interne (synthétique). & Plan dialectique construit par le candidat (thèse / antithèse / synthèse ou plan problématique). \\
\hline
\textbf{Usage du « je »} & Proscrit dans l'explication : « l'auteur montre que... », « selon Hountondji... ». & Admis avec prudence, surtout dans la synthèse et la conclusion. \\
\hline
\textbf{Référence obligatoire} & Le texte lui-même, cité et commenté ligne par ligne. & Les notions du programme et la culture philosophique personnelle. \\
\hline
\end{tabularx}
\end{center}
\legende{Tableau comparatif des deux exercices philosophiques au programme de Première technique : mémoriser ces cinq critères, c'est éviter le hors-sujet.}
\end{popfigure}

\courssub{Les trois parties de la copie et le barème}

L'épreuve sur texte n'est pas un seul exercice, mais \textbf{trois épreuves distinctes} réunies dans une même copie. Le candidat doit les traiter toutes, dans l'ordre, avec des titres visibles.

\textbf{Partie A — L'explication du texte.} Elle peut être \emph{linéaire} (on suit le texte phrase par phrase, étape par étape, en expliquant chaque argument) ou \emph{synthétique} (on dégage la structure logique globale : problème, thèse, moments de l'argumentation). Elle mobilise la compréhension et l'analyse conceptuelle.

\textbf{Partie B — Les réponses aux questions.} Deux ou trois questions précises portent sur le texte ou sur les notions de l'unité qu'il illustre (ex. : « Que reproche l'auteur à l'ethnophilosophie ? », « Distinguez tradition et traditionalisme »). Chaque réponse est courte, directe, justifiée par le texte ou par le cours.

\textbf{Partie C — L'essai personnel.} Un sujet corrélé au texte invite le candidat à prendre position de manière argumentée (ex. : « La tradition est-elle un obstacle au développement du Cameroun ? »). C'est ici, et seulement ici, que la pensée personnelle s'exprime — mais toujours de façon structurée : introduction problématisée, arguments, exemples, conclusion.

\begin{popfigure}
\begin{center}
\begin{tikzpicture}[
  grow=right,
  level 1/.style={sibling distance=3.4cm, level distance=3.6cm},
  level 2/.style={sibling distance=1.15cm, level distance=4.4cm},
  racine/.style={rectangle, rounded corners, draw=popdark, fill=popdark, text=white, font=\bfseries, align=center, inner sep=5pt},
  branche/.style={rectangle, rounded corners, draw=popblue, fill=popblueL, align=center, font=\small\bfseries, inner sep=4pt},
  feuille/.style={rectangle, rounded corners, draw=popmuted, fill=popcream, align=left, font=\scriptsize, inner sep=3pt, text width=4.1cm},
  edge from parent/.style={draw=popmuted, thick}
]
\node[racine, align=center]{Épreuve Bac\\Exercice sur texte\\(sur 20 points)}
  child { node[branche, align=center]{C. Essai personnel\\(7 à 8 pts)}
    child { node[feuille]{Temps conseillé : 45 min à 1 h} }
    child { node[feuille]{Critères : problématisation, argumentation, exemples, position personnelle} }
  }
  child { node[branche, align=center]{B. Réponses aux\\questions (4 à 5 pts)}
    child { node[feuille]{Temps conseillé : 30 à 45 min} }
    child { node[feuille]{Critères : exactitude, précision, appui sur le texte et le cours} }
  }
  child { node[branche, align=center]{A. Explication\\du texte (7 à 8 pts)}
    child { node[feuille]{Temps conseillé : 1 h à 1 h 30} }
    child { node[feuille]{Critères : fidélité à l'auteur, analyse des concepts, structure du raisonnement} }
  };
\end{tikzpicture}
\end{center}
\legende{Arborescence de l'épreuve : une racine (l'exercice sur texte), trois branches (les trois parties de la copie), et pour chacune les critères de notation et le temps conseillé pour une épreuve de 3 h à 4 h selon les séries.}
\end{popfigure}

\begin{aretenir}[Barème type et gestion du temps]
Sur 20 points, la répartition habituelle est : explication du texte \textbf{7 à 8 points}, réponses aux questions \textbf{4 à 5 points}, essai personnel \textbf{7 à 8 points} (total 20). Pour une épreuve de 3 h à 4 h selon les séries, une gestion du temps efficace est : lecture et analyse du texte au brouillon (30 min), explication (1 h à 1 h 30), questions (30 à 45 min), essai (45 min à 1 h), relecture (15 min). Ne jamais sacrifier une partie : une copie sans essai personnel perd d'office un tiers des points.
\end{aretenir}

\courssub{Première compétence à acquérir : la règle des trois lectures}

Avant même de rédiger quoi que ce soit, tout commence par la lecture. Or lire un texte philosophique n'est pas lire un roman : une seule lecture ne suffit jamais. La méthode éprouvée est celle des trois lectures successives, chacune ayant un objectif distinct.

\begin{methode}[La règle des trois lectures]
\begin{enumerate}
\item \textbf{Lecture de découverte (globale)} : lire le texte une première fois, sans stylo, pour saisir le sens général. À l'issue, formuler en une phrase : « De quoi parle ce texte ? » (le thème) puis « Que veut démontrer l'auteur ? » (la thèse). Identifier aussi l'auteur et l'œuvre si l'énoncé les donne.
\item \textbf{Lecture d'analyse (détaillée, crayon en main)} : relire lentement en soulignant les concepts clés, en entourant les articulations logiques (« donc », « car », « cependant », « en effet », « si... alors »), en numérotant les arguments dans la marge. Repérer les mots difficiles et les définir.
\item \textbf{Lecture de structuration (le plan)} : dégager les étapes du raisonnement et les transformer en plan au brouillon : quel problème est posé ? Quelle thèse est défendue ? En combien de moments l'argumentation se déploie-t-elle ? Ce plan deviendra l'ossature de la partie A.
\end{enumerate}
\end{methode}

\begin{exoresolu}[Application : premières lectures d'un extrait]
Énoncé. Voici un extrait adapté d'un texte de type bac : « On appelle ethnophilosophie la prétendue philosophie des peuples africains, présentée comme une pensée collective, anonyme et immuable. Or une philosophie n'existe que comme discours écrit, assumé par des individus qui l'argumentent et l'exposent à la discussion. La sagesse collective d'un peuple relève de la tradition, non de la philosophie. » Appliquez la règle des trois lectures : dégagez le thème, la thèse et les deux arguments.
\tcblower
\textbf{Solution détaillée.}
\begin{enumerate}
\item \emph{Lecture de découverte.} Thème : l'ethnophilosophie et la nature de la philosophie africaine. Thèse de l'auteur : l'ethnophilosophie n'est pas une philosophie au sens strict.
\item \emph{Lecture d'analyse.} Concepts clés : « ethnophilosophie » (définie dès la première phrase : pensée présentée comme collective, anonyme, immuable), « philosophie » (définie par trois critères : discours \emph{écrit}, \emph{individuel}, \emph{argumenté et discutable}). Articulation logique décisive : « Or », qui introduit l'opposition entre les deux concepts.
\item \emph{Lecture de structuration.} Plan en deux moments : (1) exposition de la thèse adverse (l'ethnophilosophie se présente comme la philosophie africaine) ; (2) réfutation en deux arguments : argument 1, la philosophie exige l'écriture et l'assomption individuelle (critères que la pensée collective ne remplit pas) ; argument 2, il faut distinguer tradition (sagesse collective) et philosophie (discours critique). Conclusion implicite : la philosophie africaine existe, mais chez des auteurs individuels, non dans la tradition anonyme.
\end{enumerate}
Ce travail de 15 minutes au brouillon fournit déjà la matière complète de la partie A de la copie.
\end{exoresolu}

\begin{attention}[Trois pièges de la lecture]
\begin{itemize}
\item \textbf{S'arrêter à la première lecture} : le sens d'un texte philosophique se découvre par couches ; qui rédige après une seule lecture paraphrase au lieu d'expliquer.
\item \textbf{Ignorer les articulations logiques} : les petits mots (« donc », « cependant », « or ») sont les charnières du raisonnement ; les manquer, c'est manquer la structure même du texte.
\item \textbf{Traduire au lieu d'expliquer} : recopier le texte « dans d'autres mots » ne vaut presque rien ; expliquer, c'est dire \emph{pourquoi} l'auteur avance cette idée, \emph{comment} elle s'articule aux autres et \emph{ce qu'elle} vise à prouver.
\end{itemize}
\end{attention}

\begin{aretenir}[L'essentiel de la section]
\begin{itemize}
\item Un texte philosophique est un écrit argumenté qui démontre une thèse par la raison, à l'aide de concepts précis.
\item L'exercice sur texte est une épreuve obligatoire du Probatoire (fin de Première) et du Baccalauréat technique, notée sur 20, évaluant compréhension, analyse, argumentation, culture et expression.
\item Il diffère radicalement de la dissertation : fidélité à l'auteur ici, création personnelle là. Confondre les deux, c'est le hors-sujet.
\item La copie comporte trois parties distinctes : explication (7--8 pts), réponses aux questions (4--5 pts), essai personnel (7--8 pts), à traiter dans un temps de 3 h à 4 h réparti méthodiquement.
\item Tout commence par la règle des trois lectures : découverte, analyse au crayon, structuration en plan.
\end{itemize}
\end{aretenir}

\begin{exercice}[Pour s'entraîner]
1. Définissez en deux phrases ce qu'est un texte philosophique et donnez un exemple de texte qui n'en est pas un, en justifiant.
2. Un élève affirme : « Expliquer un texte, c'est donner son avis sur le sujet qu'il traite. » Corrigez cette erreur en mobilisant la distinction explication / dissertation.
3. Calculez le temps maximal qu'un candidat disposant de 4 heures peut consacrer à l'essai personnel s'il respecte la répartition conseillée, et justifiez pourquoi il serait imprudent de le dépasser.
\end{exercice}

\begin{corrige}[Exercice]
1. Un texte philosophique est un écrit argumenté visant à démontrer une thèse par la raison, au moyen de concepts définis. Exemple de non-texte philosophique : un slogan publicitaire (« Votez le changement ! ») ou un cri de supporter : il affirme ou émeut sans définir ni argumenter ; il manque le problème, la thèse justifiée et la démarche rationnelle.
2. L'élève confond les deux exercices. Expliquer un texte, c'est restituer fidèlement la pensée de \emph{l'auteur} : son problème, sa thèse, ses arguments, sans y mêler son opinion. Donner son avis relève de la dissertation ou, dans l'épreuve sur texte, de la seule partie C (essai personnel). Dans la partie A, toute affirmation doit être rattachée à une ligne du texte.
3. Sur 4 h (240 min), en retirant 30 min de lectures, 1 h 30 d'explication, 45 min de questions et 15 min de relecture, il reste au maximum 60 min pour l'essai. Le dépasser serait imprudent car cela empiéterait sur la relecture (qui sauve des points d'expression) ou sur l'explication, partie la plus lourde du barème avec les questions.
\end{corrige}

\coursec{Méthodologie de l'explication de texte (Étape A)}

Dans la section précédente, nous avons montré que l'exercice sur texte philosophique n'est ni un résumé ni une dissertation déguisée : c'est un exercice de \cle{fidélité} à la pensée d'un auteur, dont l'enjeu est de faire voir au correcteur que vous avez compris \emph{comment} un texte pense, et pas seulement \emph{ce} qu'il dit. Encore faut-il savoir s'y prendre. Cette section vous donne les outils techniques de l'Étape A : l'explication proprement dite, c'est-à-dire le décorticage méthodique du texte avant toute réponse aux questions et avant tout essai personnel.

\courssub{Phase 1 : linéaire ou synthétique ? Le choix du plan d'analyse}

\begin{definition}[Explication de texte]
L'\cle{explication de texte} est une analyse ordonnée qui dégage la thèse de l'auteur, les concepts qu'il mobilise et les arguments qu'il enchaîne, en montrant la \emph{progression} de sa pensée. Elle répond à trois questions : Que dit l'auteur ? Comment le dit-il ? Pourquoi le dit-il ainsi ?
\end{definition}

Deux démarches sont possibles, et le choix n'est pas anodin :

\begin{itemize}
\item La démarche \cle{linéaire} (paragraphe par paragraphe) : on suit le texte dans l'ordre, en expliquant chaque paragraphe l'un après l'autre. Elle convient aux textes \textbf{courts} (15 à 25 lignes) et denses, où chaque paragraphe joue un rôle précis dans la démonstration.
\item La démarche \cle{synthétique} (par moments logiques) : on regroupe les paragraphes en 2 ou 3 grandes étapes de pensée (par exemple : constat, critique, proposition). Elle convient aux textes \textbf{longs} (plus de 30 lignes) ou répétitifs, où plusieurs paragraphes illustrent la même idée.
\end{itemize}

\begin{methode}[Choisir sa démarche en 3 questions]
\begin{enumerate}
\item \textbf{Compter} : le texte a-t-il 4 paragraphes ou plus ? Si oui, envisagez le regroupement par moments logiques.
\item \textbf{Comparer} : deux paragraphes disent-ils la même chose avec des exemples différents ? Si oui, fusionnez-les en un seul moment.
\item \textbf{Vérifier} : votre plan d'analyse doit toujours refléter une \emph{progression} (du problème vers la thèse, ou de la thèse vers ses conséquences), jamais une simple juxtaposition.
\end{enumerate}
\end{methode}

\begin{attention}[Erreur de jeunesse]
Au Probatoire et au Baccalauréat technique, la majorité des textes proposés compte 2 à 4 paragraphes : la démarche linéaire est alors la plus sûre. Ne choisissez la démarche synthétique que si vous savez clairement nommer chaque moment logique (« Dans un premier temps, l'auteur constate... ; dans un second temps, il critique... »). Un regroupement mal maîtrisé fait perdre plus de points qu'une analyse linéaire honnête.
\end{attention}

\courssub{Identifier le concept-clé de chaque paragraphe}

Chaque paragraphe d'un texte philosophique est construit autour d'un \cle{concept-clé} : c'est le mot-outil qui porte l'essentiel du raisonnement. Le repérer, c'est trouver la clé de voûte du paragraphe.

\begin{definition}[Concept philosophique]
Un \cle{concept philosophique} est une \emph{notion opératoire} : un outil de la pensée, défini rigoureusement par l'auteur ou par la tradition philosophique, et dont le sens ne coïncide pas toujours avec le langage courant. Par exemple, l'\emph{aliénation} chez Marx désigne la perte de maîtrise du travailleur sur le produit de son travail, et non simplement le fait d'« être fou » ; la \emph{raison} chez Kant désigne la faculté des principes, et non le simple « bon sens ».
\end{definition}

\begin{attention}[Définition philosophique, pas définition de dictionnaire]
Recopier une définition du Petit Robert est une faute méthodologique. La définition attendue est \emph{contextuelle} : elle doit dire ce que le concept signifie \textbf{dans ce texte} et \textbf{pour cet auteur}. Travaillez la définition à partir des phrases mêmes du texte (l'auteur définit souvent lui-même ses termes, parfois implicitement).
\end{attention}

Pour chaque concept repéré, vous devez produire deux choses dans votre copie :
\begin{enumerate}
\item Une \textbf{définition philosophique} brève (une à deux lignes) ;
\item Son \textbf{rôle dans l'argumentation} : sert-il à poser le problème ? à critiquer une position adverse ? à fonder la thèse ? à tirer une conséquence ?
\end{enumerate}

\begin{exemplebox}[Le concept dans son rôle]
Dans un texte où un auteur écrit que « la tradition n'est pas un musée mais un chantier », le concept-clé est \emph{tradition}. Définition philosophique : ensemble des savoirs et pratiques transmis entre générations. Rôle argumentatif : l'auteur \emph{redéfinit} le concept pour combattre la conception figée de la tradition (conception-musée) et justifier sa réappropriation critique (conception-chantier). Dire cela, c'est analyser ; répéter « la tradition c'est ce qu'on reçoit des anciens », c'est paraphraser.
\end{exemplebox}

\courssub{Repérer les connecteurs logiques : l'articulation interne du texte}

Un texte philosophique est une machine à raisonner : ses rouages visibles sont les \cle{connecteurs logiques}. Les repérer au crayon, c'est découvrir le squelette de la démonstration.

\begin{itemize}
\item \textbf{« Or »} annonce un tournant : l'auteur introduit un élément nouveau, souvent décisif, qui modifie la donne.
\item \textbf{« Mais »} marque une opposition ou une restriction : l'auteur s'écarte d'une idée qu'il vient d'évoquer.
\item \textbf{« C'est pourquoi »} introduit une conclusion : tout ce qui précède est un raisonnement qui aboutit ici.
\item \textbf{« En effet »} introduit une justification : ce qui suit prouve ou explique ce qui précède.
\end{itemize}

\begin{methode}[Fiche « Crayonnage efficace » — les 4 gestes]
\begin{enumerate}
\item \textbf{Surligner} la thèse (souvent une phrase, au début ou à la fin du texte).
\item \textbf{Encadrer} les concepts-clés (un par paragraphe en général).
\item \textbf{Flècher} les liens logiques : tracer une flèche entre la justification et ce qu'elle justifie, entre la cause et la conséquence.
\item \textbf{Numéroter} les arguments (1, 2, 3...) dans l'ordre où ils apparaissent, pour vérifier que vous n'en oubliez aucun.
\end{enumerate}
Ce crayonnage se fait \emph{sur le sujet}, en 10 minutes maximum, avant toute rédaction.
\end{methode}

\begin{exemplebox}[Boîte à outils : 20 connecteurs logiques « Bac Ready »]
\begin{center}
\begin{tabularx}{\linewidth}{|l|X|}
\hline
\rowcolor{popblueL} \textbf{Fonction} & \textbf{Connecteurs} \\
\hline
Cause & car ; parce que ; en effet ; puisque \\
\hline
Conséquence & donc ; c'est pourquoi ; ainsi ; par conséquent ; dès lors \\
\hline
Opposition & mais ; or ; cependant ; néanmoins ; en revanche \\
\hline
Concession & certes ; sans doute ; il est vrai que \\
\hline
Illustration & ainsi (exemple) ; par exemple ; c'est le cas de \\
\hline
\end{tabularx}
\end{center}
Apprenez à les classer : au brouillon, écrire « MAIS = opposition » en marge du texte vous oblige à identifier \emph{ce qui s'oppose à quoi} — c'est déjà de l'analyse.
\end{exemplebox}

\courssub{Explicite et implicite : lire entre les lignes}

\begin{definition}[Explicite et implicite]
L'\cle{explicite} est ce qui est dit directement dans le texte. L'\cle{implicite} est ce qui est supposé sans être dit : les \emph{présupposés} de l'auteur, c'est-à-dire les idées qu'il tient pour acquises et dont son raisonnement a besoin pour fonctionner.
\end{definition}

Repérer l'implicite est ce qui distingue une copie excellente d'une copie correcte. Procédez ainsi : demandez-vous « Que faut-il croire pour que cet argument soit valable ? ». Si un auteur affirme que « l'école doit former des citoyens critiques », il présuppose que l'esprit critique s'apprend et que l'école a une mission politique, pas seulement technique : ces deux présupposés ne sont pas écrits, mais ils portent tout le raisonnement.

\begin{attention}[Ne pas confondre implicite et invention]
L'implicite doit être \emph{justifié par le texte} : vous devez pouvoir montrer la phrase qui le laisse deviner. Dire « l'auteur suppose que... » exige une preuve textuelle. Inventer des intentions à l'auteur (« il pense secrètement que... ») est une faute d'interprétation.
\end{attention}

\courssub{Visualiser la progression : le croquis d'explication}

Avant de rédiger, il est utile de schématiser le texte au brouillon sous forme de flux : chaque paragraphe y apparaît avec son concept, son argument et la transition qui le relie au suivant.

\begin{popfigure}
\begin{center}
\begin{tikzpicture}[node distance=0.55cm and 0.7cm,
box/.style={draw=popblue, thick, rounded corners, fill=popblueL, text width=2.1cm, align=center, minimum height=0.9cm, font=\small},
concept/.style={draw=poporange, thick, rounded corners, fill=poporangeL, text width=1.9cm, align=center, minimum height=0.9cm, font=\small},
trans/.style={draw=popgreen, thick, fill=popgreenL, rounded corners, font=\scriptsize, align=center, text width=1.6cm},
vig/.style={draw=poppink, dashed, font=\scriptsize, text=poppink, align=center},
arr/.style={-{Stealth[length=2.5mm]}, thick, popdark}]
\node[box, align=center] (p1){\textbf{Paragraphe 1}\\ pose le problème};
\node[concept, right=of p1, align=center] (c1){Concept A\\ (définition)};
\node[box, right=of c1, fill=popgoldL, draw=popgold] (a1) {Argument 1};
\node[trans, right=of a1, align=center] (t1){Transition\\ « Or... »};
\node[box, right=of t1, align=center] (p2){\textbf{Paragraphe 2}\\ critique};
\node[concept, right=of p2] (c2) {Concept B};
\node[box, right=of c2, fill=popgoldL, draw=popgold] (a2) {Argument 2};
\draw[arr] (p1) -- (c1);
\draw[arr] (c1) -- (a1);
\draw[arr] (a1) -- (t1);
\draw[arr] (t1) -- (p2);
\draw[arr] (p2) -- (c2);
\draw[arr] (c2) -- (a2);
\node[vig, below=0.35cm of c1, align=center] (v1){Point de vigilance :\\ Définition};
\node[vig, below=0.35cm of a1, align=center] (v2){Point de vigilance :\\ Citation courte};
\node[vig, below=0.35cm of t1, align=center] (v3){Point de vigilance :\\ Analyse du lien};
\draw[poppink, dashed] (v1) -- (c1);
\draw[poppink, dashed] (v2) -- (a1);
\draw[poppink, dashed] (v3) -- (t1);
\end{tikzpicture}
\end{center}
\legende{Croquis d'explication : flux horizontal de l'analyse. Chaque paragraphe est rattaché à son concept-clé et à son argument ; les transitions (connecteurs) assurent la continuité ; les points de vigilance rappellent les trois gestes obligatoires : définir, citer, analyser.}
\end{popfigure}

Ce croquis montre aussi que les textes ne sont pas plats : la pensée de l'auteur monte en puissance. On peut représenter cette \cle{tension argumentative} comme une courbe.

\begin{popfigure}
\begin{center}
\begin{tikzpicture}
\begin{axis}[
width=0.85\linewidth, height=5.5cm,
xlabel={Paragraphes du texte},
ylabel={Intensité conceptuelle},
xmin=0.5, xmax=4.5, ymin=0, ymax=10,
xtick={1,2,3,4},
xticklabels={§1, §2, §3, §4},
ytick={2,5,8},
grid=both, grid style={popline!40},
axis lines=left]
\addplot[popcurve, domain=1:4, samples=100] {1.5 + 2.2*x - 0.15*x^2 + 0.6*sin(deg(2.5*x))};
\addplot[only marks, mark=*, mark size=2pt, poporange] coordinates {(1,3.4) (2,5.3) (3,7.4) (4,9.1)};
\node[font=\scriptsize, popdark, anchor=south] at (axis cs:4,9.1) {Thèse concluante};
\node[font=\scriptsize, popdark, anchor=south west] at (axis cs:1,3.4) {Problème posé};
\end{axis}
\end{tikzpicture}
\end{center}
\legende{Courbe simulée de la tension argumentative le long d'un texte type de 4 paragraphes : l'intensité conceptuelle croît généralement du problème initial vers la thèse finale. Votre explication doit suivre et restituer cette montée en puissance.}
\end{popfigure}

\courssub{Rédiger : les phrases de transition et la conclusion de l'explication}

L'explication rédigée n'est pas une liste d'observations : c'est un texte continu où chaque paragraphe d'analyse est relié au suivant par une \cle{phrase de transition} qui montre la progression de la pensée de l'auteur.

\begin{methode}[Construire une phrase de transition]
\begin{enumerate}
\item \textbf{Résumer} en quelques mots ce que le paragraphe analysé vient d'établir (« Après avoir montré que la tradition peut devenir un poids mort... »).
\item \textbf{Annoncer} ce que le paragraphe suivant apporte de nouveau (« ... l'auteur en tire la conséquence : seule une critique rigoureuse peut la libérer. »).
\item \textbf{Nommer} le lien logique (conséquence, approfondissement, objection réfutée) : c'est lui qui prouve que vous avez compris l'architecture du texte.
\end{enumerate}
Formules utiles : « Fort de ce constat, l'auteur... » ; « C'est précisément pourquoi... » ; « Mais cette première étape appelle un approfondissement : ... ».
\end{methode}

La \textbf{conclusion de l'explication} comporte deux mouvements obligatoires :
\begin{enumerate}
\item \textbf{Reprendre la thèse générale} de l'auteur, en une phrase reformulée avec vos mots (preuve que vous l'avez comprise) ;
\item \textbf{Évaluer la portée et les limites} de la démonstration : que prouve-t-elle réellement ? S'applique-t-elle partout (portée universelle) ou seulement dans un contexte (portée limitée) ? Un argument est-il discutable ? Cette évaluation sobre prépare l'essai personnel (Étape C) sans s'y substituer.
\end{enumerate}

\begin{exoresolu}[Analyse guidée d'un extrait inspiré de Hountondji (type Bac 2022)]
Énoncé. Un texte de 3 paragraphes, inspiré de P. Hountondji, soutient dans le premier que ce qu'on appelle « philosophie bantoue » est souvent une simple collection de croyances collectives présentées comme la pensée d'un peuple ; dans le second, que confondre philosophie et tradition collective, c'est dissoudre l'exigence de discussion rationnelle ; dans le troisième, que la philosophie africaine existe, mais comme littérature critique écrite et débattue, à l'image de toute science rigoureuse. Dégagez les concepts-clés et la thèse.
\tcblower
\textbf{Solution détaillée.}
\begin{itemize}
\item \textbf{Paragraphe 1} — Concept-clé : \emph{Tradition} (entendue comme croyance collective anonyme). Rôle : poser le problème — la « philosophie bantoue » telle que certains la décrivent n'est qu'une ethnologie déguisée. Connecteur repéré : « En effet » (justification par l'exemple des croyances).
\item \textbf{Paragraphe 2} — Concept-clé : \emph{Critique} (discussion rationnelle, possibilité du désaccord). Rôle : argument central — sans contradiction possible, il n'y a pas de philosophie. Transition : « Or » (tournant polémique contre la conception ethnophilosophique).
\item \textbf{Paragraphe 3} — Concept-clé : \emph{Scientificité} (rigueur, écriture, débat public). Rôle : renverser la conclusion pessimiste — la philosophie africaine existe, mais comme pratique critique et écrite. Connecteur : « C'est pourquoi » (conclusion).
\end{itemize}
\textbf{Thèse générale} : la philosophie, en Afrique comme ailleurs, n'est pas la tradition mais la \emph{critique rigoureuse} de la tradition. \textbf{Implicite} : l'auteur présuppose que la philosophie possède une définition universelle (exigence de rationalité) qui ne dépend pas des cultures. \textbf{Portée et limites} : la démonstration vaut pour toute prétention à identifier une « philosophie » à un folklore ; on peut toutefois discuter le présupposé d'une rationalité unique.
\end{exoresolu}

\begin{attention}[Le piège capital : paraphraser au lieu d'analyser]
\emph{Paraphraser}, c'est réécrire le texte avec d'autres mots (« l'auteur dit que la tradition n'est pas un musée »). \emph{Analyser}, c'est dire \textbf{comment} et \textbf{pourquoi} l'argument fonctionne (« l'auteur oppose deux métaphores — musée contre chantier — pour disqualifier la conception figée de la tradition et légitimer sa réappropriation critique »). Critère simple de vérification : si votre phrase peut être écrite sans avoir compris le texte, c'est de la paraphrase. Toute analyse doit contenir au moins un mot de méthode : \emph{opposer, justifier, définir, réfuter, déduire, illustrer}.
\end{attention}

\begin{aretenir}[L'essentiel de l'Étape A]
\begin{itemize}
\item Choisissez la démarche linéaire (texte court) ou synthétique (texte long) et annoncez-la.
\item Un paragraphe = un concept-clé défini philosophiquement + un rôle argumentatif identifié.
\item Crayonnez : surlignez la thèse, encadrez les concepts, flèchez les liens, numérotez les arguments.
\item Distinguez toujours l'explicite (dit) de l'implicite (présupposé, à justifier par le texte).
\item Reliez vos analyses par des phrases de transition et concluez par la thèse générale, sa portée et ses limites.
\end{itemize}
\end{aretenir}

\begin{exercice}[Entraînement]
Prenez un éditorial de journal camerounais (par exemple sur l'éducation ou la jeunesse) de 3 paragraphes. Appliquez la fiche « Crayonnage efficace » : surlignez la thèse, encadrez un concept par paragraphe, entourez tous les connecteurs logiques et classez-les selon la boîte à outils. Rédigez ensuite une phrase de transition entre le paragraphe 1 et le paragraphe 2, puis une conclusion d'explication en quatre lignes (thèse + portée).
\end{exercice}

\begin{corrige}[Exercice 1]
Éléments attendus : (1) la thèse est généralement dans le titre ou la dernière phrase de l'éditorial ; (2) les concepts-clés sont des noms abstraits répétés ou définis par l'auteur (ex. « formation », « emploi », « avenir ») ; (3) les connecteurs « mais », « car », « c'est pourquoi » doivent être classés respectivement en opposition, cause, conséquence ; (4) la phrase de transition doit résumer le §1 et annoncer le §2 en nommant le lien logique ; (5) la conclusion reprend la thèse en une phrase reformulée et indique si l'argumentation vaut universellement ou seulement dans le contexte camerounais évoqué. Toute réponse qui se contente de recopier des phrases de l'article sans mot d'analyse (\emph{opposer, justifier, conclure}) est à reprendre.
\end{corrige}

\coursec{Réponses aux questions guidées (Étape B)}

Après l'étape A, où nous avons appris à \emph{lire} le texte philosophique — repérer la thèse, dégager la structure argumentative, identifier les concepts clés —, nous abordons maintenant le cœur de l'exercice au Probatoire et au Baccalauréat technique : la \cle{réponse aux questions}. C'est ici que se gagnent ou se perdent la majorité des points. Un élève peut avoir parfaitement compris le texte et perdre la moitié de sa note parce qu'il ne sait pas \emph{rédiger} une réponse conforme aux attentes des correcteurs. Cette section vous donne les outils pour transformer votre compréhension en points.

\courssub{Typologie des questions : trois niveaux d'exigence}

Les questions posées sur un texte philosophique ne sont jamais placées au hasard. Elles suivent une progression pédagogique qui va du plus simple au plus difficile. Savoir reconnaître le type de question, c'est déjà savoir quel type de réponse construire.

\textbf{Premier niveau : les questions de compréhension.} Elles portent sur le vocabulaire, la thèse ou l'argument principal du texte. Formulations typiques : \og Que signifie l'expression \emph{homme naturel} dans ce texte ? \fg{}, \og Quelle est la thèse défendue par l'auteur ? \fg{}, \og Quel argument l'auteur oppose-t-il à la tradition ? \fg{}. La réponse se trouve \emph{dans} le texte : il faut la localiser, la citer et la reformuler avec ses propres mots. Aucune opinion personnelle n'est demandée à ce niveau.

\textbf{Deuxième niveau : les questions d'analyse.} Elles portent sur la portée d'un concept ou sur l'articulation des arguments entre eux. Formulations typiques : \og En quoi la distinction nature/culture est-elle centrale dans ce texte ? \fg{}, \og Comment l'auteur passe-t-il de la critique de l'opinion à la définition de la vérité ? \fg{}. Ici, il ne suffit plus de recopier : il faut \emph{expliquer le rôle} d'une notion dans l'économie générale du texte, montrer comment les idées s'enchaînent.

\textbf{Troisième niveau : les questions de discussion critique.} Elles invitent à prendre position : \og Partagez-vous le point de vue de l'auteur ? Justifiez votre réponse. \fg{}, \og Cette conception de la liberté vous semble-t-elle acceptable ? \fg{}. Le candidat doit exprimer un accord ou un désaccord, mais toujours de manière \emph{nuancée et argumentée}, en s'appuyant sur le texte et sur le cours.

\begin{attention}[L'erreur fatale du \og Oui / Non \fg{} sec]
Répondre \og Oui, je suis d'accord \fg{} ou \og Non, l'auteur a tort \fg{} sans nuance est une erreur éliminatoire en philosophie. La pensée philosophique exige la forme dialectique : \og Oui, mais... \fg{} ou \og Non, cependant... \fg{}. Même lorsque vous adhérez totalement à une thèse, montrez que vous en avez mesuré les limites ; même lorsque vous la rejetez, reconnaissez d'abord ce qu'elle a de solide. Un désaccord total et brutal révèle une absence de réflexion, pas une pensée forte.
\end{attention}

\courssub{La règle d'or : double justification}

Toute réponse à une question guidée repose sur une règle absolue que nous appellerons la \cle{règle de la double justification} : votre réponse doit être justifiée \textbf{par le texte} (une citation courte, entre guillemets) \textbf{et par le savoir du cours} (un auteur, une notion, un repère du programme). L'une sans l'autre ne vaut que la moitié des points.

\begin{itemize}
\item La \textbf{citation} prouve que vous avez lu le texte et que votre réponse n'est pas un hors-sujet. Elle doit être courte (une ligne maximum), exacte, et introduite (\og L'auteur affirme : \emph{...} \fg{}).
\item Le \textbf{savoir du cours} prouve que vous êtes un élève de philosophie et non un simple lecteur. Mobiliser Kant sur la raison, Descartes sur le doute, Rousseau sur la nature, montre que le texte dialogue avec votre culture philosophique.
\end{itemize}

\begin{definition}[Problématique (dans une question)]
La \cle{problématique} d'une question est la tension entre deux thèses opposées que le texte tente de résoudre. Par exemple, derrière la question \og La culture nous libère-t-elle de la nature ? \fg{} se cache la tension : d'un côté, l'homme est un être naturel soumis aux instincts ; de l'autre, il est un être culturel capable de s'élever au-dessus d'eux. Identifier cette tension, c'est comprendre ce que la question attend vraiment de vous.
\end{definition}

\courssub{Le format de la réponse idéale}

Une bonne réponse à une question guidée suit une architecture en quatre temps, que l'on peut mémoriser ainsi : \textbf{Point de vue de l'auteur} $\rightarrow$ \textbf{Preuve textuelle} $\rightarrow$ \textbf{Éclairage du cours ou exemple personnel} $\rightarrow$ \textbf{Mini-conclusion}.

\begin{enumerate}
\item \textbf{Énoncer le point de vue de l'auteur} en une ou deux phrases, sans citer encore : \og Selon l'auteur, la culture est ce qui arrache l'homme à la nature. \fg{}
\item \textbf{Apporter la preuve textuelle} : une citation courte entre guillemets, introduite et commentée.
\item \textbf{Éclairer par le cours ou par un exemple} : rapprochement avec une notion étudiée, un auteur du programme, ou un exemple concret de la vie camerounaise.
\item \textbf{Conclure brièvement} : une phrase de synthèse qui répond explicitement à la question posée.
\end{enumerate}

\begin{popfigure}
\begin{center}
\begin{tikzpicture}[node distance=0.55cm,
 bloc/.style={rectangle, rounded corners=3pt, draw=popblue, fill=popblueL, text width=4.6cm, align=center, font=\small, minimum height=1.1cm},
 fleche/.style={-{Stealth[length=3mm]}, thick, poporange}]
\node[bloc, align=center] (a){\textbf{1. Thèse de l'auteur}\\ \og Selon l'auteur... \fg{}};
\node[bloc, below=of a, fill=poporangeL, draw=poporange, align=center] (b){\textbf{2. Citation}\\ preuve textuelle courte};
\node[bloc, below=of b, fill=poppurpleL, draw=poppurple, align=center] (c){\textbf{3. Analyse conceptuelle}\\ éclairage du cours (notion, auteur)};
\node[bloc, below=of c, fill=popgreenL, draw=popgreen, align=center] (d){\textbf{4. Exemple / contre-exemple}\\ vie courante, contexte camerounais};
\node[bloc, below=of d, fill=popgoldL, draw=popgold, align=center] (e){\textbf{5. Synthèse}\\ réponse explicite à la question};
\draw[fleche] (a) -- (b);
\draw[fleche] (b) -- (c);
\draw[fleche] (c) -- (d);
\draw[fleche] (d) -- (e);
\end{tikzpicture}
\end{center}
\legende{Structure type d'une réponse idéale à une question guidée : cinq blocs enchaînés, de la thèse de l'auteur à la synthèse personnelle.}
\end{popfigure}

\begin{methode}[La technique \og C.Q.F.D. Texte \fg{}]
Pour chaque citation utilisée dans votre réponse, appliquez les quatre étapes :
\begin{enumerate}
\item \textbf{Citer} : recopier le passage entre guillemets, en indiquant si possible la ligne du texte.
\item \textbf{Qualifier} : nommer le concept philosophique en jeu (\og il s'agit ici du concept de \emph{devoir} \fg{}).
\item \textbf{Expliquer} : dégager la portée de la citation (que prouve-t-elle ? quel rôle joue-t-elle dans l'argumentation ?).
\item \textbf{Déployer} : relier explicitement cette analyse à la question posée (\og c'est donc en ce sens que l'on peut répondre que... \fg{}).
\end{enumerate}
Une citation non commentée ne rapporte aucun point : c'est votre explication qui est notée, pas votre capacité à recopier.
\end{methode}

\courssub{Gestion de l'espace et du temps}

Au Probatoire comme au Baccalauréat technique, chaque question appelle une réponse d'environ \textbf{15 à 20 lignes}. Cette contrainte est décisive et comporte deux dangers symétriques :

\begin{itemize}
\item \textbf{La réponse squelettique} (3 à 5 lignes) : elle signale au correcteur que le candidat n'a rien à dire. Même une question de simple compréhension mérite une citation, une reformulation et un éclairage.
\item \textbf{La dissertation miniature} (une page ou plus) : elle fait perdre un temps précieux pour les questions suivantes et pour l'essai personnel, et elle noye le propos. Le correcteur cherche une réponse \emph{ciblée}, pas un développement exhaustif.
\end{itemize}

La règle pratique : consacrez environ \textbf{10 à 12 minutes par question}, dont 2 minutes de repérage dans le texte, 7 minutes de rédaction et 1 minute de relecture. Si vous dépassez 20 lignes, c'est le signe que vous dérivez vers la dissertation : revenez à la question.

\begin{attention}[Le piège du \og hors texte \fg{}]
L'erreur la plus fréquente consiste à répondre uniquement avec son opinion personnelle, sans aucun ancrage dans l'extrait. Exemple : à la question \og L'auteur pense-t-il que la technique libère l'homme ? \fg{}, l'élève répond \og Oui, car grâce aux téléphones portables, les Camerounais communiquent facilement \fg{}. Cette réponse, même vraie, est \emph{hors texte} : elle ne dit rien de ce que pense l'auteur. La bonne démarche part toujours du texte (\og L'auteur affirme que... \fg{}) avant d'y adjoindre, éventuellement, un exemple personnel.
\end{attention}

\courssub{Entraînement sur corpus camerounais et africain}

L'examen porte fréquemment sur des textes de penseurs africains. Deux références majeures du programme : \textbf{Fabien Eboussi Boulaga} (philosophe camerounais, auteur de \emph{La Crise du Muntu}), qui interroge l'authenticité africaine et critique l'aliénation culturelle, et \textbf{Léopold Sédar Senghor} (poète et philosophe sénégalais), théoricien de la Négritude, pour qui \og l'émotion est nègre, comme la raison est hellène \fg{}.

\textbf{Questions types sur Eboussi Boulaga :} \og Que signifie l'expression \emph{crise du Muntu} ? \fg{} (compréhension) ; \og En quoi la critique de l'aliénation culturelle est-elle centrale dans ce texte ? \fg{} (analyse) ; \og L'authenticité africaine suppose-t-elle le rejet de l'Occident ? Discutez à partir du texte. \fg{} (discussion).

\textbf{Questions types sur Senghor :} \og Quelle conception de la connaissance l'auteur défend-il ? \fg{} (compréhension) ; \og Comment l'auteur articule-t-il émotion et raison ? \fg{} (analyse) ; \og Peut-on connaître par l'émotion ? Justifiez votre réponse. \fg{} (discussion).

\begin{exoresolu}[Question type corrigée : Rousseau et la distinction nature/culture]
\textbf{Énoncé.} À partir d'un extrait du \emph{Discours sur l'origine de l'inégalité}, on pose la question : \og En quoi la distinction nature/culture est-elle centrale dans ce texte de Rousseau ? \fg{} Rédigez une réponse de 15 à 20 lignes.
\tcblower
\textbf{Solution détaillée.}

\emph{1. Point de vue de l'auteur.} Dans cet extrait, Rousseau soutient que tout ce qui caractérise l'homme en société — l'inégalité, les passions artificielles, la vanité — n'appartient pas à sa nature mais est le produit de la culture, c'est-à-dire de la vie sociale et de l'histoire.

\emph{2. Preuve textuelle.} L'auteur écrit : \og il n'y a point de méchanceté originelle dans le cœur humain \fg{}. Cette affirmation pose que la nature humaine est un état d'innocence et d'amour de soi, et que le mal (égoïsme, domination) vient de l'institution sociale.

\emph{3. Analyse conceptuelle et éclairage du cours.} La distinction nature/culture joue ici un rôle argumentatif décisif : elle fonde le droit naturel et permet à Rousseau de critiquer la société de son temps sans condamner l'homme lui-même. Si l'inégalité est culturelle, alors elle n'est pas fatale : on peut la réformer par un contrat social légitime. Cette distinction rejoint la notion de \emph{perfectibilité}, propre à l'homme selon Rousseau : c'est parce qu'il n'est pas figé par la nature qu'il peut s'améliorer (ou se corrompre). On peut cependant en discuter la limite : la critique marxiste objectera que Rousseau néglige les déterminations économiques, et la psychanalyse freudienne rappellera que l'agressivité relève aussi de pulsions que la culture ne crée pas mais refoule.

\emph{4. Exemple.} Dans le contexte camerounais, on peut illustrer la thèse : un enfant n'est pas naturellement tribaliste ; c'est l'éducation sociale qui lui apprend à se méfier de l'autre ethnie.

\emph{5. Mini-conclusion.} La distinction nature/culture est donc centrale parce qu'elle est à la fois l'outil d'analyse du texte (elle sépare l'originaire de l'acquis) et son enjeu politique (ce qui est culturel peut être changé), même si sa rigidité peut être discutée.
\end{exoresolu}

\courssub{S'auto-évaluer : la grille du correcteur}

Pour progresser, il faut apprendre à lire sa propre copie avec les yeux du correcteur. Avant de rendre votre réponse, passez-la au crible des quatre critères suivants.

\begin{popfigure}
\begin{center}
\begin{tikzpicture}
\draw[fill=popblueL, draw=popblue] (0,0) rectangle (4.2,0.8);
\draw[fill=popblueL, draw=popblue] (4.2,0) rectangle (6.6,0.8);
\draw[fill=popblueL, draw=popblue] (6.6,0) rectangle (9.0,0.8);
\draw[fill=popblueL, draw=popblue] (9.0,0) rectangle (11.4,0.8);
\node[font=\small\bfseries] at (2.1,0.4) {Critère};
\node[font=\small\bfseries] at (5.4,0.4) {Acquis};
\node[font=\small\bfseries] at (7.8,0.4) {En cours};
\node[font=\small\bfseries] at (10.2,0.4) {Non acquis};
\foreach \y/\txt in {-0.8/{Justification textuelle}, -1.6/{Mobilisation du cours}, -2.4/{Clarté de l'expression}, -3.2/{Concision (15 à 20 lignes)}}{
  \draw (0,\y) rectangle (4.2,\y+0.8);
  \draw (4.2,\y) rectangle (6.6,\y+0.8);
  \draw (6.6,\y) rectangle (9.0,\y+0.8);
  \draw (9.0,\y) rectangle (11.4,\y+0.8);
  \node[font=\small, anchor=west] at (0.15,\y+0.4) {\txt};
};
\node[font=\small, popgreen] at (5.4,-0.4) {citation exacte};
\node[font=\small, poporange] at (7.8,-0.4) {citation vague};
\node[font=\small, poppink] at (10.2,-0.4) {aucune};
\node[font=\small, popgreen] at (5.4,-1.2) {auteur + notion};
\node[font=\small, poporange] at (7.8,-1.2) {notion seule};
\node[font=\small, poppink] at (10.2,-1.2) {hors programme};
\node[font=\small, popgreen] at (5.4,-2.0) {phrases limpides};
\node[font=\small, poporange] at (7.8,-2.0) {quelques flous};
\node[font=\small, poppink] at (10.2,-2.0) {confus};
\node[font=\small, popgreen] at (5.4,-2.8) {ciblé};
\node[font=\small, poporange] at (7.8,-2.8) {trop court/long};
\node[font=\small, poppink] at (10.2,-2.8) {dissertation};
\end{tikzpicture}
\end{center}
\legende{Grille d'auto-évaluation d'une réponse à question guidée : quatre critères, trois niveaux. Cochez mentalement chaque case avant de rendre votre copie.}
\end{popfigure}

\begin{savaistu}[Ce que cherchent les correcteurs MINESEC]
Les correcteurs du Probatoire et du Baccalauréat technique disposent d'un barème qui récompense explicitement la \emph{culture philosophique} : citer Kant sur la morale, Descartes sur le doute ou Eboussi Boulaga sur l'authenticité — à condition que ce soit \emph{à bon escient} — rapporte des points de \og culture \fg{} en plus des points de compréhension. Attention toutefois : une référence fausse ou hors de propos est sanctionnée. Mieux vaut un auteur bien maîtrisé qu'un défilé de noms mal compris.
\end{savaistu}

\begin{aretenir}[L'essentiel de l'étape B]
\begin{itemize}
\item Trois types de questions : \cle{compréhension}, \cle{analyse}, \cle{discussion critique} — chacune appelle un type de réponse différent.
\item Règle d'or : toute réponse est justifiée \textbf{par le texte} (citation courte) \textbf{et par le cours} (notion, auteur).
\item Format en quatre temps : thèse de l'auteur $\rightarrow$ preuve textuelle $\rightarrow$ éclairage du cours ou exemple $\rightarrow$ mini-conclusion.
\item 15 à 20 lignes par question, environ 10 minutes : ni squelette, ni dissertation miniature.
\item Jamais de \og oui \fg{} ou de \og non \fg{} sec : la philosophie pense en \og oui, mais... \fg{} et \og non, cependant... \fg{}.
\item Jamais d'opinion personnelle sans ancrage dans l'extrait : le texte d'abord, soi-même ensuite.
\end{itemize}
\end{aretenir}

\begin{exercice}[Entraînement guidé]
On donne l'extrait suivant, inspiré de la pensée de Fabien Eboussi Boulaga : \og Le Muntu ne se découvre authentique qu'en cessant de se définir à travers le regard de l'autre. L'aliénation commence lorsque la culture vivante se fige en folklore pour être admirée de l'extérieur. \fg{}
\begin{enumerate}
\item (Compréhension) Que signifie l'expression \og se définir à travers le regard de l'autre \fg{} ?
\item (Analyse) En quoi la distinction entre \og culture vivante \fg{} et \og folklore \fg{} est-elle centrale dans ce texte ?
\item (Discussion) L'authenticité culturelle suppose-t-elle de refuser toute influence étrangère ? Justifiez votre réponse en 15 à 20 lignes.
\end{enumerate}
\end{exercice}

\begin{corrige}[Exercice : éléments de correction]
\textbf{Question 1.} L'expression désigne le fait de construire son identité non à partir de soi-même et de sa propre expérience, mais à partir des catégories et des attentes imposées de l'extérieur, notamment par le regard colonial. Eboussi Boulaga dénonce ici une identité empruntée, qui prive le Muntu de sa capacité à se penser lui-même.

\textbf{Question 2.} La distinction est centrale car elle oppose deux états de la culture : la \emph{culture vivante}, créatrice et assumée de l'intérieur, et le \emph{folklore}, culture figée, transformée en spectacle pour le regard étranger. Cette distinction fonde le diagnostic de l'aliénation : le problème n'est pas la culture africaine elle-même, mais sa pétrification. On peut la rapprocher de la notion d'authenticité et de la critique de l'aliénation étudiées en cours.

\textbf{Question 3.} Réponse attendue en forme nuancée : \og Non, cependant... \fg{}. L'authenticité ne signifie pas l'enfermement : le texte ne condamne pas l'ouverture à l'autre, mais la \emph{définition de soi par} l'autre. Une culture peut emprunter (exemple : la musique camerounaise moderne, qui intègre des influences extérieures tout en restant ancrée dans les rythmes locaux) sans s'aliéner, à condition de rester sujet de ses choix. L'authenticité est donc une question de maîtrise de soi, non de pureté. Mini-conclusion : refuser toute influence serait d'ailleurs une nouvelle façon de se définir par rapport à l'autre — en réaction contre lui —, donc de rester aliéné.
\end{corrige}

\coursec{L'essai personnel : De la problématisation à la rédaction (Étape C)}

Après avoir répondu aux questions guidées (Étape B), qui vérifient que vous avez compris le texte et ses concepts, l'épreuve du Probatoire et du Baccalauréat technique vous demande de franchir une étape décisive : prendre position vous-même. L'essai personnel n'est pas une répétition du texte, ni une dissertation libre sans rapport avec lui. C'est un dialogue raisonné avec la thèse de l'auteur. C'est ici que le correcteur mesure votre capacité à penser de manière autonome. Cette section vous donne la méthode complète, de l'analyse du sujet jusqu'à la conclusion.

\courssub{Le lien obligatoire avec le texte}

\begin{definition}[L'essai personnel]
L'\cle{essai personnel} est un développement argumenté dans lequel le candidat prolonge, applique ou discute une thèse du texte étudié. Il ne s'agit ni de résumer le texte, ni de traiter un sujet sans rapport avec lui : le texte est le point de départ obligé de la réflexion.
\end{definition}

Imaginez un débat à la maison familiale à Yaoundé : un aîné (l'auteur du texte) a pris la parole. On ne vous demande pas de répéter ce qu'il a dit, ni de parler d'autre chose, mais de dire si vous êtes d'accord, et pourquoi. Trois rapports au texte sont possibles :

\begin{itemize}
\item \textbf{Prolonger} : vous acceptez la thèse de l'auteur et vous la développez avec vos propres arguments et exemples (par exemple, appliquer sa thèse sur la technique à la réalité du mobile money au Cameroun).
\item \textbf{Contredire} : vous montrez les limites ou les erreurs de la thèse, avec des contre-arguments solides.
\item \textbf{Nuancer} : vous acceptez partiellement la thèse et vous précisez les conditions de sa validité (c'est souvent la démarche la plus riche).
\end{itemize}

\begin{attention}[Le hors-sujet par abandon du texte]
La faute la plus grave est d'ignorer le texte : rédiger un essai général sur « la technique » ou « le travail » sans jamais faire référence à la thèse de l'auteur proposé. Le correcteur sanctionne ce hors-sujet. Règle d'or : le texte doit réapparaître au moins une fois dans l'introduction, une fois dans le développement et une fois dans la conclusion.
\end{attention}

\courssub{L'analyse du sujet : mots-clés, présupposés, champ disciplinaire}

Avant toute rédaction, il faut démonter le sujet comme un mécanicien démonte un moteur avant de le réparer. Trois opérations :

\begin{enumerate}
\item \textbf{Dégager les mots-clés} : ce sont les termes porteurs du sens (ex. dans « La technique libère-t-elle l'homme ? » : \emph{technique}, \emph{libère}, \emph{homme}). Chacun doit être défini au brouillon.
\item \textbf{Repérer les présupposés} : tout sujet suppose quelque chose sans le dire. « La technique libère-t-elle l'homme ? » présuppose que l'homme pourrait être prisonnier de quelque chose, et que la technique a un pouvoir sur lui. Interroger ces présupposés, c'est déjà philosopher.
\item \textbf{Identifier le champ disciplinaire} : le sujet relève-t-il de l'éthique (le bien, le devoir), de la philosophie politique (l'État, la justice), de l'épistémologie (la science, la vérité), de la philosophie de l'action (le travail, la technique) ? Cette identification oriente le choix des auteurs et des exemples.
\end{enumerate}

\begin{definition}[Problématiser]
\cle{Problématiser}, ce n'est pas définir les mots du sujet : c'est faire surgir le \emph{conflit} d'interprétations que le sujet rend possible. Un sujet philosophique est toujours traversé par une tension : d'un côté une réponse semble s'imposer (la thèse), de l'autre une réponse opposée est tout aussi défendable (l'antithèse). La problématique est la question qui exprime cette tension.
\end{definition}

Par exemple, pour « La technique libère-t-elle l'homme ? » : d'un côté, la technique semble libérer (elle réduit l'effort, allège les tâches, rapproche les hommes) ; de l'autre, elle semble asservir (dépendance au téléphone, chômage technique, contrôle social). La problématique tensionnelle sera donc : \emph{la technique est-elle un instrument de liberté ou une nouvelle forme de servitude ?}

\courssub{Le choix du plan}

Une fois la tension dégagée, il faut organiser la réflexion. Trois plans sont admis au Probatoire et au Baccalauréat technique :

\begin{center}
\begin{tabularx}{\linewidth}{|l|X|X|}
\hline
\rowcolor{popblueL} \textbf{Plan} & \textbf{Structure} & \textbf{Quand l'utiliser ?} \\
\hline
\textbf{Dialectique} & I. Thèse — II. Antithèse — III. Synthèse (dépassement) & Sujet sous forme de question fermée opposant deux réponses (« Faut-il... ? », « Est-ce... ? ») \\
\hline
\textbf{Analytique} & I. Causes ou genèse — II. Structure ou fonctionnement — III. Conséquences ou finalité & Sujet portant sur un phénomène à expliquer \\
\hline
\textbf{Thématique} & I, II, III : trois aspects complémentaires de la notion & Sujet invitant à décrire les dimensions d'un concept \\
\hline
\end{tabularx}
\end{center}

Le \cle{plan dialectique} est le plus fréquent et le plus sûr pour l'essai personnel : il épouse naturellement la tension de la problématique. La synthèse n'est pas un compromis mou (« les deux ont raison »), mais un dépassement : elle montre sous quelles conditions la thèse et l'antithèse peuvent être conciliées.

\begin{attention}[La sanction du « plan catalogue »]
Le \cle{plan catalogue} consiste à aligner des auteurs sans lien logique : « Descartes dit ceci, Kant dit cela, moi je dis ceci. » C'est une juxtaposition, pas une argumentation. Le correcteur exige l'\textbf{articulation} : chaque partie doit naître de la précédente (la partie II répond à une difficulté de la partie I ; la partie III dépasse l'opposition des deux premières). Interrogez-vous : si j'inverse mes parties, le raisonnement tient-il encore ? Si oui, c'est un catalogue.
\end{attention}

\courssub{La rédaction : introduction, développement, conclusion}

\textbf{L'introduction (méthode AMPEC).} Elle se construit en cinq mouvements, comme un entonnoir qui va du large vers le précis :

\begin{itemize}
\item \textbf{A — Accroche} : un fait, une situation concrète ou une constatation qui amène naturellement au sujet (ex. : « À Douala comme à Maroua, le téléphone portable a transformé le commerce, les études et même les relations familiales. »).
\item \textbf{M — Mots-clés définis} : définir brièvement les termes essentiels du sujet.
\item \textbf{P — Problématique} : formuler la question tensionnelle.
\item \textbf{E — Énoncé de la thèse} : indiquer la position que vous allez défendre.
\item \textbf{C — Chemin (plan annoncé)} : annoncer les étapes du développement.
\end{itemize}

\textbf{Le développement.} Il comporte deux ou trois grandes parties. Chaque partie obéit à une structure fixe : une \emph{idée force} (énoncée dans la première phrase), des \emph{arguments} (raisonnements), des \emph{exemples} (de préférence tirés de la vie camerounaise et africaine), un appui sur un \emph{auteur} cité brièvement, et une \emph{transition} qui prépare la partie suivante.

\textbf{La conclusion.} Elle comporte deux temps : le \emph{bilan synthétique} (réponse claire à la problématique, en rappelant le chemin parcouru) et l'\emph{ouverture} (une question nouvelle : autre auteur, autre domaine, ou fait d'actualité).

\begin{popfigure}
\begin{center}
\begin{tikzpicture}[node distance=0.35cm,
 bloc/.style={draw=popblue, fill=popblueL, rounded corners=2pt, minimum width=4.6cm, minimum height=0.75cm, font=\small, align=center},
 trans/.style={draw=poporange, fill=poporangeL, rounded corners=2pt, minimum width=3.4cm, minimum height=0.5cm, font=\footnotesize\itshape, align=center},
 entonnoir/.style={draw=popgreen, fill=popgreenL, trapezium, trapezium angle=70, minimum width=5.2cm, minimum height=0.9cm, font=\small, align=center},
 entonnoirinv/.style={draw=popgreen, fill=popgreenL, trapezium, trapezium angle=70, shape border rotate=180, minimum width=5.2cm, minimum height=0.9cm, font=\small, align=center},
 ouvert/.style={draw=poppurple, fill=poppurpleL, rounded corners=2pt, minimum width=4.2cm, minimum height=0.55cm, font=\footnotesize, align=center},
 fleche/.style={-{Stealth[length=2.5mm]}, thick, popdark}]
\node[entonnoir, align=center] (intro){Introduction (AMPEC)\\ \footnotesize du général vers le sujet};
\node[bloc, below=of intro, align=center] (p1){Partie I : Thèse\\ \footnotesize idée force + arguments + exemples};
\node[trans, below=of p1] (t1) {Transition : limite de la thèse};
\node[bloc, below=of t1, align=center] (p2){Partie II : Antithèse\\ \footnotesize idée force + arguments + exemples};
\node[trans, below=of p2] (t2) {Transition : vers le dépassement};
\node[bloc, below=of t2, align=center] (p3){Partie III : Synthèse\\ \footnotesize conditions de conciliation};
\node[entonnoirinv, below=of p3, align=center] (concl){Conclusion : bilan synthétique\\ \footnotesize réponse à la problématique};
\node[ouvert, below=of concl] (ouv) {Ouverture : autre auteur, autre domaine, actualité};
\draw[fleche] (intro) -- (p1);
\draw[fleche] (p1) -- (t1);
\draw[fleche] (t1) -- (p2);
\draw[fleche] (p2) -- (t2);
\draw[fleche] (t2) -- (p3);
\draw[fleche] (p3) -- (concl);
\draw[fleche] (concl) -- (ouv);
\end{tikzpicture}
\end{center}
\legende{Architecture de la dissertation type Bac : l'introduction et la conclusion fonctionnent comme des entonnoirs (du large au précis, puis du précis au large), les parties sont des blocs argumentés reliés par des transitions.}
\end{popfigure}

\courssub{Du brouillon à la copie : gestion du temps}

\begin{methode}[Le brouillon en 3 temps]
\begin{enumerate}
\item \textbf{Analyse du sujet (10 min)} : souligner les mots-clés, dégager les présupposés, formuler la problématique tensionnelle, choisir le plan.
\item \textbf{Plan détaillé (30 min)} : pour chaque partie, noter au brouillon l'idée force, deux arguments, un auteur, un exemple concret, et la phrase de transition.
\item \textbf{Rédaction propre (1 h 20)} : rédiger l'introduction au brouillon puis au propre, puis le développement et la conclusion directement, en surveillant l'orthographe et en relisant les cinq dernières minutes.
\end{enumerate}
\end{methode}

\begin{popfigure}
\begin{center}
\begin{tikzpicture}[mindmap, grow cyclic,
 every node/.style={concept, font=\small},
 concept color=popblue!60,
 level 1/.append style={level distance=3.1cm, sibling angle=72, font=\footnotesize},
 level 2/.append style={level distance=2.1cm, sibling angle=45, font=\scriptsize}]
\node[concept color=popblue!70] {La technique libère-t-elle l'homme ?}
 child[concept color=popgreen!60] { node {Définitions}
   child { node {Technique : ensemble d'outils et de procédés} }
   child { node {Liberté : pouvoir d'agir par soi-même} } }
 child[concept color=poporange!60] { node {Auteurs}
   child { node {Marx : machine et aliénation} }
   child { node {Heidegger : la technique comme arraisonnement} }
   child { node {Ellul : système technicien} } }
 child[concept color=poppurple!60] { node {Arguments POUR}
   child { node {Allège l'effort physique} }
   child { node {Accès au savoir, communication} } }
 child[concept color=poppink!60] { node {Arguments CONTRE}
   child { node {Dépendance, addiction aux écrans} }
   child { node {Chômage technique} } }
 child[concept color=popgold!70] { node {Exemples Cameroun}
   child { node {Mobile money : autonomie financière} }
   child { node {Motoculteurs et agriculture modernisée} } };
\end{tikzpicture}
\end{center}
\legende{Carte mentale de brainstorming pour le sujet « La technique libère-t-elle l'homme ? » : cinq branches (définitions, auteurs, arguments pour, arguments contre, exemples camerounais) alimentent le plan détaillé au brouillon.}
\end{popfigure}

\begin{exemplebox}[Un plan complet sur un sujet de type Bac]
\textbf{Sujet (Bac technique 2024) : « Le travail est-il une aliénation ? »}
\begin{itemize}
\item \textbf{I. Le travail comme aliénation} : dans la production en série, l'ouvrier ne se reconnaît pas dans le produit de son travail (Marx, \emph{Manuscrits de 1844}) ; exemple : les ouvriers des zones industrielles de Douala, payés à la tâche, sans maîtrise du produit fini.
\item \textbf{II. Le travail comme humanisation} : par le travail, l'homme transforme la nature et se transforme lui-même ; il se reconnaît dans son œuvre (Hegel, dialectique du maître et de l'esclave) ; exemple : l'artisan menuisier de Foumban ou le potier, fiers de leur savoir-faire.
\item \textbf{III. Dépassement : les conditions d'un travail libérateur} : le travail n'aliène pas par nature, mais selon ses conditions ; un cadre juridique et social juste le rend émancipateur ; exemple : le Code du travail camerounais, les coopératives agricoles de l'Ouest où le producteur reste maître de sa récolte.
\end{itemize}
\end{exemplebox}

\begin{exoresolu}[Rédiger une introduction complète (AMPEC)]
\textbf{Sujet : « La technique libère-t-elle l'homme ? »} Rédiger l'introduction de l'essai personnel.
\tcblower
\textbf{Solution détaillée.}
\emph{Accroche} : « En quelques années, le téléphone mobile a transformé la vie quotidienne des Camerounais : on paie ses factures, on envoie de l'argent, on suit des cours à distance. »
\emph{Mots-clés définis} : « La technique désigne l'ensemble des outils et des procédés inventés par l'homme pour agir sur la nature ; la liberté est le pouvoir d'agir selon sa propre volonté. »
\emph{Problématique} : « Mais cet outil qui semble nous rendre plus autonomes ne nous rend-il pas aussi dépendants ? La technique est-elle un instrument de liberté ou une nouvelle forme de servitude ? »
\emph{Énoncé de thèse} : « Nous défendrons l'idée que la technique libère l'homme à condition qu'il en demeure le maître. »
\emph{Plan annoncé} : « Nous montrerons d'abord que la technique libère l'homme des contraintes naturelles, puis qu'elle peut aussi l'asservir, avant d'examiner les conditions d'un usage libérateur de la technique. »
\end{exoresolu}

\begin{savaistu}[L'ouverture valorisée]
Au Baccalauréat technique, l'ouverture finale sur un fait de société camerounais récent est très appréciée des correcteurs, à condition qu'elle soit philosophiquement pertinente : une grève d'enseignants (question de la justice sociale), le débat sur la sous-traitance dans les grandes entreprises (question de l'aliénation par le travail), l'essor des réseaux sociaux chez les jeunes (question de la liberté et de la technique). Une ouverture n'est pas une conclusion bis : c'est une question nouvelle qui montre que la réflexion continue.
\end{savaistu}

\begin{exercice}[À vous de jouer]
À partir du texte étudié en classe sur le travail, on propose le sujet d'essai : « Travailler plus pour vivre mieux ? »
\begin{enumerate}
\item Dégager les mots-clés et deux présupposés du sujet.
\item Formuler une problématique tensionnelle.
\item Construire un plan dialectique en trois parties, avec pour chaque partie une idée force, un auteur et un exemple camerounais.
\end{enumerate}
\end{exercice}

\begin{corrige}[Exercice]
\begin{enumerate}
\item Mots-clés : \emph{travailler}, \emph{plus}, \emph{vivre mieux}. Présupposés : (a) le bien-être dépendrait de la quantité de travail ; (b) « vivre mieux » se réduirait à des biens matériels.
\item Problématique : « L'augmentation du travail est-elle la condition d'une vie meilleure, ou la vie bonne exige-t-elle au contraire de limiter le travail ? »
\item Plan dialectique : \textbf{I.} Travailler plus permet de vivre mieux : le travail procure revenu, dignité et reconnaissance sociale (Hegel ; exemple : l'entrepreneur de Bafoussam qui développe son atelier). \textbf{II.} Mais l'excès de travail détruit la vie : fatigue, perte du lien familial, aliénation (Marx ; exemple : les journaliers des plantations de la côte, épuisés pour un salaire minimal). \textbf{III.} Dépassement : vivre mieux suppose un travail juste, mesuré et doté de sens, encadré par le droit (le Code du travail camerounais fixant la durée légale du travail) : c'est la qualité du travail, non sa quantité, qui humanise.
\end{enumerate}
\end{corrige}

\begin{aretenir}[L'essentiel de l'étape C]
L'essai personnel part toujours du texte (prolonger, contredire ou nuancer sa thèse). La démarche : analyser le sujet (mots-clés, présupposés, champ), problématiser en faisant surgir le conflit thèse/antithèse, choisir un plan (dialectique le plus souvent), rédiger selon AMPEC pour l'introduction, des blocs « idée force + arguments + exemples + transition » pour le développement, et une conclusion en deux temps (bilan + ouverture). Le temps se gère en trois temps : 10 min d'analyse, 30 min de plan détaillé, 1 h 20 de rédaction.
\end{aretenir}

\coursec{TP Intégré : Simulation Bac sur corpus « Philosophie et Développement au Cameroun »}

Dans la section précédente, nous avons appris à construire un essai personnel : problématiser un sujet, bâtir un plan dialectique, rédiger une introduction et une conclusion. Mais savoir faire et \emph{faire dans les conditions réelles} sont deux choses différentes. Au Probatoire comme au Baccalauréat technique, vous ne disposerez que de quelques heures, sans cours sous les yeux, face à un texte inconnu. Cette séance de travaux pratiques est donc une \cle{simulation complète} : elle met en œuvre, dans l'ordre réel de l'épreuve, l'explication de texte (étape A), les réponses aux questions (étape B) et l'essai personnel (étape C). L'objectif n'est pas seulement de produire une copie, mais d'apprendre à \emph{s'évaluer soi-même} comme le ferait un correcteur officiel.

\courssub{Le support de la simulation : un sujet format Bac officiel}

\textbf{Présentation du corpus}\par
Le texte support est un extrait de \emph{Consciencisme} de Kwame Nkrumah (1964), œuvre au programme MINESEC dans l'unité « Philosophie et Développement ». Nkrumah, premier président du Ghana, y défend la thèse suivante : le développement de l'Afrique exige une philosophie propre, synthèse des traditions africaines, de l'héritage islamique et de la colonisation européenne, capable de guider l'action politique et économique. Ce texte permet de mobiliser les notions de l'unité : \cle{développement}, \cle{tradition}, \cle{modernité}, \cle{aliénation culturelle}, \cle{conscience}.

\begin{exemplebox}[Le sujet complet tel qu'il sera projeté]
Le document projeté au vidéoprojecteur reproduit la mise en page officielle d'un sujet du Baccalauréat technique :
\begin{itemize}
\item \textbf{En-tête} : « République du Cameroun — Paix, Travail, Patrie / MINESEC — Inspection de Philosophie / Baccalauréat de l'Enseignement Technique ».
\item \textbf{Mentions réglementaires} : Durée : 3 heures ; Coefficient : 2 ; Aucun document autorisé.
\item \textbf{Texte} : extrait de Nkrumah, \emph{Consciencisme} (15 à 20 lignes), suivi de la consigne : « Expliquez le texte suivant ».
\item \textbf{Question 1} : « Dégagez la thèse de l'auteur et montrez comment il la justifie. »
\item \textbf{Question 2} : « Le développement d'un pays comme le Cameroun suppose-t-il nécessairement la rupture avec les traditions ? »
\item \textbf{Sujet d'essai} : « La philosophie peut-elle contribuer au développement du Cameroun ? »
\end{itemize}
\end{exemplebox}

\begin{attention}[Gestion du temps réel de l'épreuve]
Le piège classique est de tout miser sur l'essai et de bâcler l'explication de texte, qui rapporte pourtant souvent la moitié des points. Répartition conseillée pour une épreuve de 2 h 30 à 3 h : \textbf{30 min} pour l'explication de texte, \textbf{30 min} pour les réponses aux questions, \textbf{1 h 30} pour l'essai. Réservez impérativement \textbf{15 min de relecture} finale : orthographe, cohérence entre les parties, vérification que la thèse annoncée en introduction est bien celle défendue en conclusion. Une copie brillante mais inachevée est toujours moins bien notée qu'une copie honnête et complète.
\end{attention}

\courssub{Activité 1 (en groupe) : Crayonnage collaboratif au tableau}

\textbf{Intuition d'abord.} Avant d'écrire quoi que ce soit, un bon lecteur \emph{voit} le texte : il repère d'un coup d'œil où est la thèse, quels mots portent l'argumentation, comment les étapes s'enchaînent. C'est exactement ce que fait un correcteur quand il lit votre copie. Le crayonnage collectif entraîne ce regard.

\begin{methode}[Protocole du crayonnage collectif (15 min)]
\begin{enumerate}
\item Le texte est projeté au vidéoprojecteur. Lecture à voix haute par un élève, sans interruption.
\item Un premier groupe vient \textbf{souligner en bleu} la ou les phrases qui expriment la \cle{thèse} de l'auteur.
\item Un deuxième groupe \textbf{encadre en rouge} les \cle{concepts} philosophiques (ici : conscience, tradition, révolution, développement).
\item Un troisième groupe \textbf{numérote les étapes} de l'argumentation dans la marge et relie les \cle{connecteurs logiques} (car, donc, cependant, en effet).
\item La classe valide ou corrige chaque annotation ; le professeur n'intervient qu'en dernier recours.
\end{enumerate}
\end{methode}

À l'issue de cette activité, la structure du texte de Nkrumah apparaît clairement : (1) constat — l'Afrique est marquée par trois héritages culturels concurrents ; (2) problème — sans philosophie unificatrice, ces héritages restent sources de confusion ; (3) thèse — seule une philosophie de la conscience, le « consciencisme », peut fonder un développement authentique. Retenir ce réflexe : \textbf{toute explication de texte commence par une structure, pas par une paraphrase}.

\courssub{Activité 2 (individuel, 20 min) : Réponses aux questions et correction par les pairs}

Chaque élève rédige, sur copie double et sans document, ses réponses aux deux questions. Rappel de la règle d'or de l'étape B : une bonne réponse est \emph{courte, structurée et justifiée} — elle cite le texte pour la question 1, elle mobilise les notions de l'unité et un exemple camerounais pour la question 2.

Puis les copies sont échangées en croisé (rangée A vers rangée B) et évaluées à l'aide de la grille ci-dessous, inspirée des grilles officielles de l'Inspection.

\begin{popfigure}
\begin{center}
\begin{tikzpicture}[scale=1]
% En-têtes colonnes (niveaux)
\fill[popblue] (2.6,4.4) rectangle (5.2,5.2); \node[white,font=\small\bfseries] at (3.9,4.8) {Insuffisant};
\fill[poporange] (5.2,4.4) rectangle (7.8,5.2); \node[white,font=\small\bfseries] at (6.5,4.8) {Passable};
\fill[popgreen] (7.8,4.4) rectangle (10.4,5.2); \node[white,font=\small\bfseries] at (9.1,4.8) {Bien};
\fill[poppurple] (10.4,4.4) rectangle (13.0,5.2); \node[white,font=\small\bfseries] at (11.7,4.8) {Excellent};
\node[font=\small\bfseries] at (3.9,5.5) {0--2 pts};
\node[font=\small\bfseries] at (6.5,5.5) {3--4 pts};
\node[font=\small\bfseries] at (9.1,5.5) {5--6 pts};
\node[font=\small\bfseries] at (11.7,5.5) {7--8 pts};
% Lignes de critères
\foreach \y/\crit/\a/\b/\c/\d in {
  3.3/{Compréhension}/{Hors sujet ou paraphrase}/{Thèse repérée mais imprécise}/{Thèse et structure dégagées}/{Sens global restitué avec finesse},
  2.2/{Analyse}/{Aucun concept identifié}/{Concepts cités sans définition}/{Concepts définis et articulés}/{Analyse des étapes et des enjeux},
  1.1/{Argumentation}/{Simple opinion}/{Arguments généraux}/{Arguments + exemples pertinents}/{Dialectique (thèse/antithèse) maîtrisée},
  0.0/{Expression}/{Langue très défaillante}/{Fautes fréquentes}/{Langue claire et correcte}/{Style précis, vocabulaire philosophique}}{
  \fill[popblueL] (0,\y-0.5) rectangle (2.6,\y+0.5);
  \node[font=\small\bfseries,anchor=west] at (0.1,\y) {\crit};
  \node[font=\scriptsize,anchor=west,text width=2.4cm] at (2.7,\y) {\a};
  \node[font=\scriptsize,anchor=west,text width=2.4cm] at (5.3,\y) {\b};
  \node[font=\scriptsize,anchor=west,text width=2.4cm] at (7.9,\y) {\c};
  \node[font=\scriptsize,anchor=west,text width=2.4cm] at (10.5,\y) {\d};
  \draw[popline] (0,\y-0.5) rectangle (13.0,\y+0.5);
  \draw[popline] (2.6,\y-0.5) -- (2.6,\y+0.5);
  \draw[popline] (5.2,\y-0.5) -- (5.2,\y+0.5);
  \draw[popline] (7.8,\y-0.5) -- (7.8,\y+0.5);
  \draw[popline] (10.4,\y-0.5) -- (10.4,\y+0.5);
}
\draw[popline] (0,3.8) -- (13.0,3.8);
\draw[popline] (2.6,4.4) -- (2.6,5.2);
\draw[popline] (5.2,4.4) -- (5.2,5.2);
\draw[popline] (7.8,4.4) -- (7.8,5.2);
\draw[popline] (10.4,4.4) -- (10.4,5.2);
\end{tikzpicture}
\end{center}
\legende{Grille d'évaluation croisée des réponses aux questions (barème sur 8 points par critère, inspirée des grilles de l'Inspection Pédagogique de Philosophie). Chaque évaluateur coche une case par ligne et additionne.}
\end{popfigure}

\begin{methode}[Protocole « Correction par les pairs »]
\begin{enumerate}
\item \textbf{Lis} attentivement la copie de ton voisin, en silence, une première fois sans rien écrire.
\item \textbf{Coche} la grille d'évaluation, critère par critère, en justifiant mentalement chaque choix.
\item \textbf{Écris} au bas de la copie \emph{un point fort} (formulé positivement : « ta thèse est clairement dégagée ») et \emph{un axe de progrès} (formulé comme un conseil : « cite le texte ligne 4 pour appuyer ton analyse »).
\item \textbf{Restitution orale} : chaque évaluateur explique sa notation à l'auteur de la copie en 1 minute ; l'auteur peut se défendre en citant sa propre copie.
\end{enumerate}
\end{methode}

\begin{exoresolu}[Question 2 : corrigé type]
\textbf{Énoncé.} « Le développement d'un pays comme le Cameroun suppose-t-il nécessairement la rupture avec les traditions ? » Rédigez une réponse de 15 à 20 lignes.
\tcblower
\textbf{Solution détaillée.} Une bonne réponse procède en trois temps. \emph{Premier temps} : définir les termes. Le développement désigne la transformation durable des conditions matérielles et humaines d'une société ; la tradition est l'ensemble des pratiques et valeurs transmises de génération en génération. \emph{Deuxième temps} : confronter deux thèses. On peut soutenir, avec les partisans de la modernité, que certaines traditions freinent le progrès (par exemple des pratiques successorales qui excluent les femmes de l'héritage foncier dans certaines régions du Cameroun). Mais on peut soutenir, avec Nkrumah, qu'un développement qui nie les traditions est un développement aliéné : il copie des modèles étrangers sans racines, comme ces projets industriels importés clé en main qui échouent faute d'appropriation locale. \emph{Troisième temps} : dépasser l'alternative. Le développement authentique n'exige ni la conservation intégrale ni la rupture totale, mais un \emph{tri critique} : conserver ce qui fonde la solidarité (la tontine, le travail communautaire, réemployés avec succès dans l'épargne et le financement de projets locaux) et abandonner ce qui contredit la dignité humaine. \textbf{Conclusion} : la tradition n'est pas un obstacle en soi ; c'est l'absence de réflexion philosophique sur la tradition qui est un obstacle.
\end{exoresolu}

\courssub{Activité 3 (en groupe) : Brainstorming et construction collective du plan de l'essai}

Le sujet d'essai — « La philosophie peut-elle contribuer au développement du Cameroun ? » — est écrit au tableau. Chaque groupe dispose de 5 minutes pour proposer : une \cle{problématique}, deux arguments opposés, et un exemple camerounais. Le professeur note toutes les propositions sans les juger (règle du brainstorming : aucune idée n'est rejetée pendant la production), puis la classe procède au tri.

La construction collective aboutit typiquement à ce plan dialectique, validé par le professeur :
\begin{itemize}
\item \textbf{Thèse} : la philosophie semble inutile au développement, affaire d'ingénieurs, d'économistes et de financements (le bitume des routes ne se discute pas, il se construit).
\item \textbf{Antithèse} : mais aucun développement durable n'est possible sans réflexion sur ses fins — développer \emph{pour quoi}, \emph{pour qui} ? Nkrumah montre que les politiques sans philosophie conduisent à l'imitation et à l'échec.
\item \textbf{Synthèse} : la philosophie contribue au développement non comme technique mais comme \emph{conscience critique} : elle forme des citoyens capables de choisir leur avenir, condition de tout progrès authentique au Cameroun.
\end{itemize}

\begin{attention}[Piège fréquent lors du brainstorming]
Ne confondez pas \emph{accumuler des idées} et \emph{construire un plan}. Une liste de dix exemples n'est pas un plan : le plan exige une \textbf{progression} (chaque partie répond à une objection de la précédente) et une \textbf{hiérarchie} (la synthèse doit dépasser l'opposition, pas simplement la répéter). Vérifiez toujours que vos trois parties peuvent être reliées par « mais » et « donc ».
\end{attention}

\courssub{Activité 4 (individuel, 45 min) : Rédaction en conditions réelles}

Chaque élève rédige, sur copie double, l'\textbf{introduction complète} (accroche, définitions, problématique, annonce du plan) et la \textbf{première partie du développement} (la thèse), dans le silence et la posture d'une salle d'examen : téléphone éteint, aucun échange, gestion autonome du temps. Le professeur annonce le temps restant à mi-parcours, comme le surveillant au Bac.

Les copies sont ensuite ramassées et corrigées selon le protocole « Correcteur Bac » : notation sur 20, annotations en marge (jamais de simple « faux » : chaque annotation indique \emph{pourquoi} et \emph{comment améliorer}), et appréciation globale distinguant le fond (logique, pertinence) de la forme (langue, présentation).

\begin{savaistu}[Le « hack » légal du candidat camerounais]
L'Inspection Générale Pédagogique en charge de la Philosophie publie chaque année un \textbf{rapport de jury} qui analyse les erreurs majoritaires des candidats : hors-sujets récurrents, notions confondues, citations apprises par cœur et plaquées hors de propos. Lire ce rapport, c'est connaître à l'avance les pièges que le correcteur a sous les yeux quand il lit votre copie. C'est une stratégie parfaitement légale — et étonnamment peu utilisée. Demandez-le à votre professeur ou consultez les archives du MINESEC.
\end{savaistu}

\courssub{Retour méta-cognitif : analyser des copies types}

La dernière étape est la plus formatrice : la classe étudie trois \textbf{copies anonymisées} fournies par l'Inspection Pédagogique — une copie faible, une copie moyenne, une copie excellente — portant sur le même sujet. Il ne s'agit plus de philosopher, mais de \emph{penser sur sa propre pensée} : c'est la \cle{méta-cognition}.

\begin{exemplebox}[Ce que révèle la comparaison des trois copies]
\begin{itemize}
\item \textbf{Copie faible (6/20)} : paraphrase le texte ligne par ligne, aucune thèse dégagée, l'essai raconte « le développement au Cameroun » sans répondre à la question « la philosophie peut-elle... ». Leçon : le hors-sujet partiel est le premier tueur de notes.
\item \textbf{Copie moyenne (11/20)} : thèse correcte mais non justifiée, exemples génériques (« par exemple en Afrique »), plan annoncé mais non tenu. Leçon : la différence entre 10 et 14 se joue sur la \emph{justification}, pas sur les idées.
\item \textbf{Copie excellente (16/20)} : problématique explicite dès l'introduction, chaque argument appuyé sur le texte ou sur une notion du programme, exemples précis et camerounais, conclusion qui répond réellement à la question posée. Leçon : l'excellence est une affaire de \emph{méthode}, non de génie.
\end{itemize}
\end{exemplebox}

\begin{exemplebox}[Résultats chiffrés d'une expérimentation réelle]
Lors de la simulation de 2023 au Lycée Technique de Douala, la moyenne de la classe de Première technique est passée de \textbf{8,5/20} à \textbf{12,4/20} après seulement trois séances de TP utilisant la correction par les pairs et l'analyse de copies types — soit un gain de près de 4 points, supérieur à celui obtenu par trois séances de cours magistral supplémentaires. S'évaluer soi-même fait plus progresser qu'écouter.
\end{exemplebox}

\begin{aretenir}[Kit de survie « Bac Philo »]
\begin{enumerate}
\item \textbf{Fiche mémo des concepts clés} de l'unité : développement, tradition, modernité, conscience, aliénation — chacun avec une définition et un exemple camerounais.
\item \textbf{Liste de connecteurs} : d'abord / ensuite / enfin (plan) ; cependant / néanmoins (objection) ; en effet / car (justification) ; ainsi / par conséquent (conclusion).
\item \textbf{Structure de l'introduction} : accroche $\rightarrow$ définition des termes $\rightarrow$ problématique $\rightarrow$ annonce du plan. \textbf{Structure de la conclusion} : réponse claire à la problématique $\rightarrow$ bilan du parcours $\rightarrow$ ouverture courte.
\item \textbf{Gestion du temps} : 30 min explication, 30 min questions, 1 h 30 essai, 15 min relecture.
\item \textbf{Entraînement} : refaire au moins cinq sujets probables en conditions réelles avant l'examen, dont un corpus « Philosophie et Développement au Cameroun ».
\end{enumerate}
\end{aretenir}

\begin{exercice}[Pour aller plus loin]
Reprenez votre copie de l'activité 4 corrigée. (a) Réécrivez uniquement votre introduction en intégrant les annotations du correcteur. (b) À l'aide de la grille 4$\times$4, auto-évaluez votre nouvelle version et calculez votre progression en points. (c) Formulez par écrit, en trois lignes, l'erreur que vous vous engagez à ne plus commettre le jour de l'examen.
\end{exercice}

\begin{corrige}[Exercice — TP Intégré]
Éléments de correction : (a) La nouvelle introduction doit contenir les quatre moments obligatoires ; vérifiez que la problématique reprend les termes exacts du sujet (« la philosophie », « contribuer », « développement du Cameroun ») et que l'annonce du plan est explicite (« Nous verrons d'abord que..., ensuite que..., enfin que... »). (b) Comparez les totaux obtenus critère par critère : la progression doit apparaître principalement sur « Argumentation » et « Expression ». (c) L'engagement doit être précis et vérifiable (« je citerai le texte au moins deux fois par partie ») et non vague (« je ferai plus attention »). Le professeur collecte ces engagements et les redistribue la veille de l'examen blanc suivant.
\end{corrige}

\coursec{Erreurs fréquentes, Remédiation et Excellence}

Le TP intégré de la section précédente vous a placés en situation réelle d'examen : corpus sur « Philosophie et Développement au Cameroun », questions guidées, essai personnel. Vous avez sans doute constaté une chose : connaître la méthode ne suffit pas. Entre la méthode comprise et la copie réussie, il y a un fossé, et ce fossé est presque toujours le même pour tous les candidats. Cette dernière section a donc un objectif précis : identifier les erreurs qui coûtent des points, proposer des exercices de remédiation ciblés, et définir ce qui distingue une copie correcte (10--12/20) d'une copie d'excellence (16/20 et plus). Car au Probatoire comme au Baccalauréat technique, l'explication de texte ne se gagne pas seulement en cours : elle se gagne dans les détails de la copie.

\courssub{Le Top 5 des erreurs au Bac technique}

Avant de corriger, il faut diagnostiquer. L'observation des copies de l'enseignement technique fait ressortir cinq erreurs récurrentes, classées ici par ordre de gravité.

\textbf{Erreur n°1 : le hors-sujet (thème contre texte).} Le candidat lit le texte, repère un mot — par exemple « travail » — et rédige tout ce qu'il sait sur le travail, sans jamais revenir au texte. Or l'exercice porte sur \emph{ce} texte, avec \emph{sa} thèse particulière. Dire « le travail est une aliénation » alors que l'auteur défend que le travail humanise l'homme, c'est un hors-sujet, même si la phrase est philosophiquement intéressante.

\textbf{Erreur n°2 : la paraphrase.} Le candidat recopie le texte en changeant quelques mots : « L'auteur dit que le développement ne se réduit pas à la croissance économique » devient « L'auteur affirme que le développement n'est pas seulement la croissance économique ». Il n'y a aucune analyse : pas de concept dégagé, pas d'argument explicité, pas d'enjeu formulé. Le correcteur sanctionne : « Vous répétez, vous n'expliquez pas. »

\textbf{Erreur n°3 : le « je pense que » sans argument.} Dans l'essai personnel, affirmer « je pense que la tradition est un obstacle » sans donner ni raison, ni exemple, ni référence, c'est donner une opinion, pas une position philosophique. En philosophie, toute thèse doit être \emph{justifiée}.

\textbf{Erreur n°4 : le plan invisible.} Les parties existent dans la tête du candidat mais pas dans la copie : pas d'alinéa, pas de transitions, pas de connecteurs. Le correcteur ne peut pas suivre la progression et conclut à l'absence de structure.

\textbf{Erreur n°5 : le français approximatif.} Accords négligés, syntaxe bancale, vocabulaire imprécis (« la chose », « le truc », « ça montre »). La langue est le véhicule de la pensée : une pensée floue s'exprime dans une langue floue, et inversement.

\begin{attention}[Le piège du sens commun]
Utiliser « liberté », « justice », « travail », « État » au sens courant (celui des conversations quotidiennes) au lieu du sens philosophique défini dans le cours est une faute majeure. Dire « la liberté, c'est faire ce qu'on veut » relève du sens commun ; dire « la liberté est l'autodétermination rationnelle de la volonté » relève du concept philosophique. Sanction du correcteur : « Confusion des genres ». Avant d'utiliser un grand mot, demandez-vous toujours : \emph{quelle est sa définition philosophique précise ?}
\end{attention}

\begin{definition}[La rigueur philosophique]
La \cle{rigueur philosophique} est l'adéquation exacte entre trois éléments : le \cle{concept} utilisé (défini précisément), l'\cle{argument} déployé (logiquement enchaîné) et la \cle{thèse} défendue (clairement formulée). Il n'y a pas de « flou artistique » en philosophie : chaque affirmation doit pouvoir être rattachée à un raisonnement et à une définition.
\end{definition}

\courssub{Remédiation ciblée : trois ateliers}

Chaque erreur appelle un exercice correcteur spécifique. Voici trois ateliers à pratiquer régulièrement, seul ou en classe.

\begin{methode}[Atelier 1 : la reformulation analytique]
Objectif : transformer une phrase de l'auteur en véritable analyse. Procédez en quatre étapes :
\begin{enumerate}
\item \textbf{Repérer} la phrase-clé du texte et la citer entre guillemets.
\item \textbf{Dégager le concept} : quel concept philosophique cette phrase met-elle en jeu ?
\item \textbf{Expliciter l'argument} : quel raisonnement l'auteur mène-t-il (cause, conséquence, opposition) ?
\item \textbf{Formuler l'enjeu} : pourquoi cette thèse est-elle importante ? Que remet-elle en question ?
\end{enumerate}
Formule magique : « En affirmant que \og ... \fg, l'auteur entend montrer que ..., car ..., ce qui signifie que ... ».
\end{methode}

\begin{exoresolu}[De la paraphrase à l'analyse]
Énoncé. Texte : « Le développement d'un pays ne se mesure pas à ses richesses, mais à la qualité de vie de son peuple. » Transformez la paraphrase suivante en analyse : « L'auteur dit que le développement n'est pas la richesse mais la qualité de vie. »
\tcblower
Solution. La paraphrase se contente de répéter. L'analyse procède ainsi : « En affirmant que le développement \og ne se mesure pas à ses richesses \fg, l'auteur oppose deux concepts : la \emph{croissance économique} (quantité de biens produits) et le \emph{développement humain} (qualité de vie : santé, éducation, liberté). Son argument repose sur une distinction entre la fin et les moyens : la richesse n'est qu'un moyen, la fin véritable est le bien-être du peuple. L'enjeu est politique : un pays comme le Cameroun peut connaître une croissance du PIB sans que la vie quotidienne des populations de l'Extrême-Nord ou de l'Est s'améliore réellement. » On est passé du « quoi » (paraphrase) au « comment » et au « pourquoi » (analyse).
\end{exoresolu}

\begin{methode}[Atelier 2 : la chasse au connecteur]
Objectif : rendre le plan visible. Prenez un de vos anciens devoirs et soulignez tous les connecteurs logiques. Si vous en trouvez moins de huit, votre structure est invisible. Intégrez alors les connecteurs « Premium » : \emph{d'abord / ensuite / enfin} (progression), \emph{cependant / toutefois / néanmoins} (opposition), \emph{en effet / car / puisque} (justification), \emph{par conséquent / ainsi / c'est pourquoi} (conclusion), \emph{non seulement... mais encore} (amplification).
\end{methode}

\begin{methode}[Atelier 3 : la rédaction de transitions types]
Objectif : articuler les parties de l'essai personnel. Une bonne transition fait deux choses : elle \textbf{bilanne} la partie qui se termine et elle \textbf{annonce} la partie suivante, souvent sous forme de question. Modèles à mémoriser :
\begin{enumerate}
\item « Si donc [thèse de la partie 1], toutefois [objection qui ouvre la partie 2] : [question]. »
\item « Nous avons établi que [bilan]. Reste à savoir si [nouveau problème]. »
\item « Mais cette solution est-elle suffisante ? Ne faudrait-il pas plutôt [partie 3] ? »
\end{enumerate}
Entraînement : rédigez trois transitions différentes pour un même plan en deux parties sur « La tradition est-elle un obstacle au développement ? ».
\end{methode}

\courssub{Les critères de l'excellence : viser plus de 16/20}

Qu'est-ce qui sépare une copie de 12 d'une copie de 18 ? Quatre critères, que le correcteur repère immédiatement.

\begin{aretenir}[Les quatre critères d'excellence]
\begin{enumerate}
\item \textbf{Originalité de la problématique} : ne pas poser une question banale (« Qu'est-ce que la liberté ? ») mais une question qui naît d'une tension réelle (« Si le développement exige l'obéissance à des règles collectives, est-il compatible avec la liberté individuelle ? »).
\item \textbf{Pertinence des exemples} : croiser un exemple classique (Platon, Descartes) et un exemple local (le débat sur les langues nationales à l'école, le rôle des tontines dans l'économie camerounaise, la question de l'accès à l'eau potable).
\item \textbf{Maîtrise conceptuelle précise} : chaque concept est défini avant d'être utilisé, et jamais confondu avec un autre (ne pas mélanger « croissance » et « développement », « loi » et « règle »).
\item \textbf{Style élégant} : phrases courtes, vocabulaire précis, rythme, aucune lourdeur. L'élégance est la clarté devenue agréable.
\end{enumerate}
\end{aretenir}

\begin{popfigure}
\begin{center}
\begin{tikzpicture}[
  box/.style={draw=popblue, fill=popblueL, rounded corners=3pt, text width=4.6cm, minimum height=1.1cm, font=\small, align=center},
  note/.style={font=\footnotesize\itshape, text=popdark, text width=4.8cm, align=left}
]
\node[box] (intro) at (0,0) {\textbf{Introduction AMPEC} : Amorce, Mise au point, Problématique, Enjeu, Chemin};
\node[note, right=0.4cm of intro] {Accroche concrète + question problématisée, pas de définition de dictionnaire.};
\node[box, fill=popgreenL, draw=popgreen] (trans) at (0,-2) {\textbf{Transition type} : « Si la tradition assure la cohésion, toutefois fige-t-elle l'innovation ? »};
\node[note, right=0.4cm of trans] {Bilan + question : le plan devient visible et logique.};
\node[box, fill=poporangeL, draw=poporange] (ex) at (0,-4) {\textbf{Exemple local intégré} : tontines, forage villageois, débat sur les langues};
\node[note, right=0.4cm of ex] {L'exemple illustre le concept, il ne le remplace pas.};
\node[box, fill=poppurpleL, draw=poppurple] (ouv) at (0,-6) {\textbf{Ouverture conceptuelle} : du développement économique vers le développement humain};
\node[note, right=0.4cm of ouv] {La conclusion élargit sans quitter le sujet.};
\draw[->, thick, popblue] (intro) -- (trans);
\draw[->, thick, popgreen] (trans) -- (ex);
\draw[->, thick, poporange] (ex) -- (ouv);
\end{tikzpicture}
\end{center}
\legende{Anatomie d'une copie 18/20 : quatre marqueurs repérables par le correcteur en moins de deux minutes de lecture.}
\end{popfigure}

\begin{exemplebox}[Objectif « zéro faute de français »]
Au Probatoire et au Baccalauréat, environ \textbf{1 point sur 20} est fréquemment dédié à la qualité de la langue : accords, ponctuation, vocabulaire précis. Un point, c'est souvent la différence entre 9 et 10, entre l'ajournement et l'admission. Technique de relecture efficace : \textbf{relire sa copie à l'envers}, de la dernière phrase vers la première. Le cerveau, ne pouvant plus suivre le sens, se concentre uniquement sur la forme et repère les fautes d'accord et les répétitions. Réservez systématiquement les dix dernières minutes de l'épreuve à cette relecture.
\end{exemplebox}

\begin{popfigure}
\begin{center}
\begin{tikzpicture}
\begin{axis}[
  ybar, bar width=0.8cm,
  width=0.85\linewidth, height=6cm,
  ymin=0, ymax=35,
  ylabel={Points perdus (en \%)},
  symbolic x coords={Hors-sujet, Paraphrase, Expression, Structure, Culture},
  xtick=data,
  x tick label style={font=\small, rotate=15, anchor=east},
  nodes near coords={\pgfmathprintnumber{\pgfplotspointmeta}\,\%},
  nodes near coords style={font=\small\bfseries},
  ymajorgrids, grid style={popline!50},
  axis line style={popdark}
]
\addplot[fill=popblue, draw=popdark] coordinates {(Hors-sujet,30) (Paraphrase,25) (Expression,20) (Structure,15) (Culture,10)};
\end{axis}
\end{tikzpicture}
\end{center}
\legende{Répartition estimée des points perdus à l'explication de texte (enquête sur copies d'examen, 2023). Le hors-sujet et la paraphrase représentent à eux seuls 55 \% des pertes : travailler l'analyse est donc la priorité absolue.}
\end{popfigure}

\courssub{Gérer le stress et organiser son temps}

L'excellence ne dépend pas que de l'intelligence : elle dépend de la maîtrise de soi le jour de l'épreuve.

\textbf{La respiration.} Avant d'ouvrir le sujet, pratiquez la respiration 4--6 : inspirez 4 secondes par le nez, expirez 6 secondes par la bouche, cinq fois de suite. Cette technique ralentit le rythme cardiaque et libère l'accès à la mémoire. Un candidat stressé oublie ce qu'il sait ; un candidat calme mobilise tout son cours.

\textbf{La lecture active.} Ne lisez jamais un texte philosophique les mains vides : crayon en main, soulignez la thèse, encadrez les connecteurs de l'auteur (« car », « donc », « mais »), notez en marge les concepts. La lecture active transforme un texte intimidant en un document de travail.

\textbf{Le planning horaire.} Écrivez votre plan de bataille sur la table (ou le brouillon) dès le début de l'épreuve et tenez-vous-y :

\begin{center}
\begin{tabularx}{\linewidth}{|l|X|}
\hline
\rowcolor{popblueL} \textbf{Durée} & \textbf{Tâche} \\
\hline
20 min & Lecture active du texte, repérage de la thèse et des concepts \\
\hline
15 min & Réponses aux questions au brouillon (plan détaillé) \\
\hline
45 min & Rédaction des réponses aux questions \\
\hline
15 min & Problématisation et plan de l'essai personnel \\
\hline
60 min & Rédaction de l'essai personnel \\
\hline
10 min & Relecture à l'envers (traque des fautes) \\
\hline
\end{tabularx}
\end{center}

\courssub{Ressources numériques et fiche de révision}

\begin{savaistu}[La dissertation des séries techniques]
L'exercice philosophique en série technique (F, STI et filières assimilées) attend davantage d'exemples \emph{concrets et techniques} qu'en série littéraire. Adaptez votre banque d'exemples à votre spécialité : l'intelligence artificielle et l'automatisation pour parler du travail, les énergies renouvelables (barrages, solaire) pour parler du développement, la gestion de projet et le travail d'équipe pour parler de la vie en société. Un exemple tiré de votre spécialité technique, bien analysé, impressionne toujours favorablement le correcteur.
\end{savaistu}

Pour vous entraîner, exploitez les ressources numériques mises à disposition dans le cadre du programme MINESEC : la plateforme \textbf{« École Numérique »} propose une banque de sujets corrigés d'explications de texte ; la chaîne YouTube \textbf{« Philo Cameroun »} offre des capsules de méthode courtes (5 à 10 minutes) à revoir la veille de l'épreuve. L'entraînement sur sujets réels reste la préparation la plus efficace.

\begin{methode}[La fiche de révision d'une page]
Concevez une fiche recto-verso, à relire la veille de l'examen :
\begin{itemize}
\item \textbf{Recto} : la méthode en 3 étapes (analyse du texte, réponses guidées, essai personnel) + le plan type de l'essai (introduction problématisée, deux ou trois parties avec transitions, conclusion ouverte).
\item \textbf{Verso} : 10 concepts « Joker » définis en une ligne (liberté, devoir, État, justice, travail, technique, culture, bonheur, vérité, développement) + 5 connecteurs « Premium » + la check-list de relecture (thèse citée ? plan visible ? exemples locaux ? relecture à l'envers faite ?).
\end{itemize}
\end{methode}

\courssub{Projection post-Bac : des compétences pour la vie}

Terminons par une question légitime : à quoi servira cet exercice après le Baccalauréat ? La réponse est : partout. Les trois compétences travaillées dans l'explication de texte — \textbf{l'analyse critique} (comprendre un document complexe et en dégager l'essentiel), \textbf{l'argumentation} (défendre une position avec des raisons), \textbf{la rédaction structurée} (présenter clairement une pensée) — sont des atouts majeurs :

\begin{itemize}
\item pour les \textbf{concours administratifs} (ENAM, IRIC, concours de la fonction publique), où la dissertation et la note de synthèse sont décisives ;
\item pour les \textbf{Grandes Écoles} et l'enseignement supérieur, où tout rapport, tout mémoire exige de problématiser et d'argumenter ;
\item pour l'\textbf{entrepreneuriat}, où convaincre un investisseur, rédiger un plan d'affaires ou défendre un projet devant une banque suppose exactement les mêmes compétences : une problématique claire, des arguments solides, des exemples concrets, un style convaincant.
\end{itemize}

\begin{aretenir}[L'essentiel de la section]
Les cinq erreurs à éliminer : hors-sujet, paraphrase, opinion sans argument, plan invisible, français approximatif. Les trois ateliers de remédiation : reformulation analytique, chasse au connecteur, transitions types. Les quatre critères d'excellence : problématique originale, exemples croisés (locaux et classiques), concepts précis, style élégant. Enfin, la rigueur philosophique — adéquation du concept, de l'argument et de la thèse — est une compétence qui vous suivra bien au-delà de l'examen : c'est la marque d'un esprit formé.
\end{aretenir}

\newpage

\coursec{Activité d'intégration}

\coursec{TP Intégré : Simulation Bac sur corpus « Philosophie et Développement au Cameroun »}

\courssub{Objectifs de l'activité}

Cette activité d'intégration, intitulée \cle{Le Grand Oral Philo Tech}, place l'élève dans une situation authentique d'évaluation orale inspirée des entretiens post-Baccalauréat (concours d'entrée dans les grandes écoles, ENAM, ENS, écoles de génie, entretiens de bourse). Elle vise à mobiliser simultanément les trois savoirs de la leçon :

\begin{itemize}
\item les \cle{généralités} sur la nature et les enjeux de l'exercice sur texte philosophique ;
\item les \cle{réponses aux questions} guidées, y compris face à une question inattendue ;
\item l'\cle{essai personnel}, de la problématisation à la défense d'une thèse.
\end{itemize}

Au terme de l'activité, l'élève doit être capable de :
\begin{enumerate}
\item dégager oralement, en un temps limité, la thèse et l'argumentation d'un texte philosophique ;
\item répondre avec rigueur et réactivité à une question imprévue ;
\item construire et défendre le plan d'un essai personnel en justifiant sa démarche ;
\item évaluer la prestation d'un pair à l'aide d'une grille critériée, ce qui développe le recul critique.
\end{enumerate}

\courssub{Situation}

\textbf{Contexte.} Le lycée technique de Nkongsamba organise, une semaine avant les épreuves du Probatoire technique, sa traditionnelle \emph{Journée de l'éloquence philosophique}. Le proviseur a invité deux anciens élèves, aujourd'hui étudiants à l'ENSP de Yaoundé, pour témoigner : \og au concours, on ne vous demande pas de réciter, on vous demande de \emph{penser à voix haute} à partir d'un texte \fg. Chaque élève de la classe de Première technique doit donc s'entraîner à défendre une thèse face à un jury.

\textbf{Règle du jeu.} L'élève-candidat tire au sort une enveloppe contenant :
\begin{itemize}
\item un \cle{texte court} (10 à 15 lignes) d'un auteur africain au programme ;
\item un \cle{sujet d'essai} corrélé au texte.
\end{itemize}

\textbf{Enveloppe n\textsuperscript{o} 3 (celle qui servira de corrigé commenté).}

\emph{Texte} (inspiré de la pensée d'Eboussi Boulaga, philosophe camerounais) : \og La liberté ne se reçoit pas comme une aumône ; elle se conquiert dans la pratique. Un peuple qui attend que son développement lui soit offert de l'extérieur renonce déjà à être l'auteur de sa propre histoire. Penser, pour l'Africain, c'est d'abord se déprendre des modèles imposés, interroger sa situation concrète, et produire à partir d'elle les moyens de son propre avenir. \fg

\emph{Sujet d'essai corrélé} : Le développement du Cameroun peut-il être pensé sans une liberté de pensée effective de ses citoyens ?

\textbf{Données numériques de l'épreuve.}
\begin{center}
\begin{tabularx}{\linewidth}{|l|X|}\hline
\rowcolor{popblueL} \textbf{Phase} & \textbf{Durée et contenu} \\ \hline
Préparation & 20 minutes, fiche méthode de la leçon autorisée, brouillon permis \\ \hline
Étape 1 : Explication synthétique & 3 minutes : thèse, arguments, structure du texte \\ \hline
Étape 2 : Question piège du jury & 2 minutes : une question inattendue posée par le jury \\ \hline
Étape 3 : Essai personnel & 5 minutes : problématique, plan annoncé, défense de la thèse \\ \hline
Débriefing collectif & 10 minutes : \og Qu'est-ce qui a différencié le candidat convaincant ? \fg \\ \hline
\end{tabularx}
\end{center}

\textbf{Composition du jury} : deux élèves tirés au sort et le professeur. Le jury note sur 20 selon la grille simplifiée suivante :

\begin{center}
\begin{tabularx}{\linewidth}{|l|X|c|}\hline
\rowcolor{popblueL} \textbf{Critère} & \textbf{Indicateurs observables} & \textbf{Barème} \\ \hline
Clarté & Expression structurée, vocabulaire précis, temps respecté & /5 \\ \hline
Rigueur & Concepts définis, arguments ordonnés, pas de contradiction & /5 \\ \hline
Réactivité & Réponse pertinente à la question piège, sans panique ni digression & /5 \\ \hline
Ancrage dans le texte & Citations courtes, références exactes à l'auteur & /5 \\ \hline
\end{tabularx}
\end{center}

\courssub{Consignes}

\textbf{Travail du candidat (pendant les 20 minutes de préparation).}

\begin{enumerate}
\item Lis le texte deux fois. Souligne la thèse de l'auteur et formule-la en une seule phrase, sans recopier le texte.
\item Repère les deux ou trois arguments qui soutiennent cette thèse et classe-les du plus faible au plus fort.
\item Dégage l'enjeu philosophique du texte : de quelle notion du programme relève-t-il ? Quel problème concret camerounais illustre-t-il ?
\item Pour le sujet d'essai : transforme la question du sujet en problématique, puis construis un plan en deux ou trois parties (thèse, antithèse, dépassement, ou plan thématique progressif).
\item Prévois une objection possible du jury et prépare ta réponse.
\end{enumerate}

\textbf{Travail du jury (pendant la prestation).}

\begin{enumerate}
\setcounter{enumi}{5}
\item Note chaque critère au fur et à mesure, en justifiant chaque note par une observation précise (exemple : \og cite le texte deux fois \fg, \og ne définit pas le mot liberté \fg).
\item Formule la question piège : elle doit porter sur le texte ou sur le sujet, être courte, et exiger une prise de position (exemples autorisés : \og La liberté se conquiert, dites-vous : le colonisé était-il donc responsable de sa servitude ? \fg ; \og Un peuple peut-il se développer sans modèle extérieur du tout ? \fg).
\item Lors du débriefing, proposez collectivement un conseil d'amélioration formulé de manière positive.
\end{enumerate}

\textbf{Travail de la classe (observateurs).}

\begin{enumerate}
\setcounter{enumi}{8}
\item Remplis la fiche d'écoute : quelle est la thèse défendue ? Quel est le moment le plus convaincant de la prestation, et pourquoi ?
\end{enumerate}

\courssub{Ressources à mobiliser}

\begin{aretenir}[Ce qu'il faut avoir en tête avant de passer devant le jury]
\begin{itemize}
\item \textbf{Explication de texte} : thèse (ce que l'auteur soutient) $\rightarrow$ arguments (comment il le prouve) $\rightarrow$ enjeu (pourquoi cela importe). Toujours citer brièvement le texte pour ancrer son propos.
\item \textbf{Réponse à une question} : écouter jusqu'au bout, reformuler si besoin, répondre d'abord par \emph{oui}, \emph{non} ou \emph{cela dépend}, puis justifier en deux arguments maximum.
\item \textbf{Essai personnel} : problématique claire, plan annoncé, thèse personnelle assumée et défendue par des arguments et des exemples concrets (vie courante, réalité camerounaise).
\item \textbf{Posture orale} : regarder le jury, parler posément, gérer son temps (un minuteur est visible).
\end{itemize}
\end{aretenir}

\courssub{Démarche et corrigé commenté}

\begin{exoresolu}[Prestation modèle sur l'enveloppe n\textsuperscript{o} 3]
\textbf{Énoncé.} À partir du texte inspiré d'Eboussi Boulaga et du sujet \og Le développement du Cameroun peut-il être pensé sans une liberté de pensée effective de ses citoyens ? \fg, produire la prestation complète attendue en 10 minutes.
\tcblower
\textbf{Étape 1 — Explication synthétique (3 min).}

\og Ce texte défend une thèse : la liberté et le développement ne se reçoivent pas, ils se conquièrent par la pratique et par la pensée. L'auteur avance trois arguments. D'abord, un argument d'expérience : attendre le développement de l'extérieur, c'est \og renoncer à être l'auteur de sa propre histoire \fg ; l'aide extérieure ne remplace pas l'initiative. Ensuite, un argument de définition : penser, pour l'Africain, ce n'est pas copier des modèles, c'est \og interroger sa situation concrète \fg. Enfin, un argument de finalité : le but de cette pensée est de \og produire les moyens de son propre avenir \fg. L'enjeu est donc philosophique et politique : il s'agit de l'autonomie de la pensée comme condition du développement. \fg

\emph{Commentaire du jury} : la thèse est formulée en une phrase, les trois arguments sont ordonnés, deux citations courtes ancrent le propos dans le texte. Critère \og ancrage \fg : 5/5.

\textbf{Étape 2 — Réponse à la question piège (2 min).}

Question du jury : \og Si la liberté se conquiert, les peuples colonisés étaient-ils responsables de leur propre servitude ? \fg

Réponse modèle : \og Non, et le texte ne dit pas cela. Dire que la liberté se conquiert ne signifie pas que l'opprimé est coupable de son oppression : la colonisation était une domination imposée par la force. Mais cela signifie que la libération ne peut pas être offerte par l'oppresseur : elle suppose l'action et la pensée des opprimés eux-mêmes. L'auteur parle d'ailleurs au présent, à l'Africain d'aujourd'hui : c'est à nous, maintenant, de ne pas reproduire une dépendance mentale. \fg

\emph{Commentaire} : réponse qui commence par une prise de position nette (\og Non \fg), distingue deux sens de \og conquérir \fg, et revient au texte. Réactivité : 5/5.

\textbf{Étape 3 — Essai personnel : plan et défense de la thèse (5 min).}

Problématique : \og Le développement est-il d'abord une affaire de moyens matériels (routes, usines, capitaux) ou d'abord une affaire de liberté de pensée ? \fg

Plan annoncé : (1) Le développement exige des moyens matériels et techniques — sans eux, aucune pensée ne suffit ; (2) mais ces moyens restent stériles ou imposés si les citoyens ne pensent pas librement leurs besoins — d'où les échecs des projets copiés de l'extérieur ; (3) donc la liberté de pensée est la condition qui rend les moyens efficaces : former des techniciens capables d'innover, comme le fait l'enseignement technique au Cameroun, c'est déjà développer le pays.

Thèse défendue : \og Je soutiens que le développement du Cameroun ne peut pas être pensé sans liberté de pensée effective, parce qu'un technicien qui ne fait qu'appliquer des modes d'emploi étrangers ne résoudra jamais les problèmes propres à son pays. L'exemple des jeunes qui fabriquent localement des machines agricoles adaptées à nos sols montre que penser par soi-même produit du développement réel. \fg

\emph{Commentaire} : problématique explicite, plan dialectique en trois temps, thèse personnelle assumée avec un exemple camerounais concret. Note globale : 18/20 (un point retiré pour un léger dépassement de temps, un point pour une définition du mot \og développement \fg restée implicite).
\end{exoresolu}

\courssub{Bilan de l'activité}

Cette activité a permis de mettre en évidence plusieurs acquis essentiels de la leçon.

\begin{itemize}
\item \textbf{L'unité des trois exercices.} Expliquer un texte, répondre à une question et rédiger un essai ne sont pas trois tâches séparées : l'explication fournit la matière, la question entraîne la réactivité, l'essai couronne le tout par une prise de position personnelle. Le Grand Oral les enchaîne exactement comme l'examen écrit.
\item \textbf{Le temps comme contrainte méthodologique.} En 20 minutes de préparation, seule une démarche ordonnée (thèse $\rightarrow$ arguments $\rightarrow$ problématique $\rightarrow$ plan) permet d'être prêt ; c'est la même discipline qu'à l'épreuve écrite du Probatoire et du Baccalauréat technique.
\item \textbf{Ce qui différencie le candidat convaincant.} Le débriefing collectif fait émerger les critères décisifs : une thèse formulée en une phrase, des citations courtes et exactes, une réponse franche aux questions pièges, et un exemple concret tiré de la réalité camerounaise. À l'inverse, la récitation de notions sans lien avec le texte, les digressions et l'absence de prise de position sont sanctionnées par la grille.
\item \textbf{L'évaluation par les pairs.} En notant un camarade, chaque élève intériorise les critères de réussite et apprend à lire sa propre copie avec le regard du correcteur.
\end{itemize}

\begin{attention}[Pour réussir le jour de l'oral]
Ne jamais paniquer devant une question piège : elle ne vise pas à humilier mais à vérifier que la thèse est comprise et assumée. La bonne stratégie est toujours la même : prise de position nette, distinction conceptuelle si nécessaire, retour au texte. Un silence de trois secondes pour organiser sa pensée vaut mieux qu'une réponse confuse de deux minutes.
\end{attention}

\newpage

\coursec{Méthodes et exercices résolus}

\coursec{L'exercice sur texte philosophique}

\begin{aretenir}[Généralités : nature et enjeux de l'exercice]
L'exercice sur texte philosophique est l'une des deux formes de l'épreuve de philosophie au Probatoire et au Baccalauréat technique (l'autre étant la dissertation). Il évalue la capacité du candidat à :
\begin{itemize}
\item \textbf{Comprendre} un texte philosophique : en dégager la thèse, la structure argumentative, les concepts clés et les exemples.
\item \textbf{Analyser} le texte en répondant à des questions guidées qui exigent explication, justification par citation et commentaire.
\item \textbf{Problématiser} : à partir de la tension interne au texte, construire une question philosophique et y répondre par un essai personnel structuré (thèse, antithèse, synthèse).
\end{itemize}
\textbf{Durée} : 3 heures. \textbf{Notation} : généralement 6 à 8 points pour les questions guidées, 12 à 14 points pour l'essai personnel. La rigueur méthodologique (respect des consignes, citation commentée, plan dialectique) prime sur l'érudition.
\end{aretenir}

\courssub{Appliquer les notions de l'unité à des situations concrètes de la vie courante}

\begin{methode}[Analyser un texte philosophique : la lecture méthodique]
Avant toute réponse, il faut comprendre le texte. Voici la démarche en 5 étapes :
\begin{enumerate}
\item \textbf{Identifier le texte} : relever l'auteur, l'œuvre, le thème général (ex. : la technique, la liberté, l'État). Cela montre au correcteur que vous situez le passage.
\item \textbf{Dégager la thèse} : demandez-vous « que veut démontrer l'auteur ? ». Formulez-la en une phrase simple : « L'auteur soutient que... ».
\item \textbf{Repérer la structure} : soulignez les connecteurs logiques (\emph{car}, \emph{donc}, \emph{mais}, \emph{en effet}, \emph{cependant}). Ils découpent le texte en moments de l'argumentation.
\item \textbf{Définir les termes-clés} : chaque mot important (ex. : \cle{technique}, \cle{développement}, \cle{liberté}) doit être défini avec vos propres mots, sans dictionnaire.
\item \textbf{Relever les arguments et exemples} : notez dans la marge chaque argument et l'exemple qui l'illustre, en indiquant sa fonction (prouver, illustrer, réfuter).
\end{enumerate}
\textbf{Réflexe Bac} : ne jamais commencer à rédiger avant d'avoir formulé la thèse par écrit au brouillon.
\end{methode}

\begin{exoresolu}[Dégager la thèse d'un texte]
\textbf{Énoncé.} Texte : « L'homme n'est pas seulement un être de nature ; par la technique, il transforme le monde et se transforme lui-même. Mais cette puissance comporte un risque : l'outil peut asservir celui qui l'a fabriqué. Le développement n'a donc de sens que s'il demeure au service de l'homme. »
\begin{enumerate}
\item Dégagez la thèse du texte.
\item Relevez deux connecteurs logiques et précisez leur fonction.
\end{enumerate}
\tcblower
\textbf{Solution détaillée.}
\begin{enumerate}
\item \textbf{Recherche de la thèse.} Je m'interroge : que veut démontrer l'auteur ? Le texte affirme d'abord que la technique est un pouvoir de transformation (de la nature et de l'homme), puis signale un danger (l'asservissement par l'outil), et conclut par une condition. La phrase conclusive, introduite par « donc », porte l'idée principale.\\
\textbf{Thèse} : l'auteur soutient que la technique n'est un véritable facteur de développement que si elle reste au service de l'homme et non l'inverse.
\item \textbf{Connecteurs logiques.}
\begin{itemize}
\item « \textbf{Mais} » : connecteur d'opposition ; il introduit la restriction ou le danger (l'outil qui asservit), c'est-à-dire le moment critique de l'argumentation.
\item « \textbf{Donc} » : connecteur de conséquence ; il introduit la conclusion du raisonnement, c'est-à-dire la thèse elle-même.
\end{itemize}
\end{enumerate}
\textbf{Conclusion.} Le texte suit un mouvement en trois temps : constat (la technique transforme), objection (elle peut asservir), conclusion (elle doit rester au service de l'homme). Identifier ce mouvement, c'est avoir compris le texte.
\end{exoresolu}

\begin{methode}[Répondre aux questions guidées : la règle Q-R-J]
Chaque question portée sur le texte exige une réponse en trois temps :
\begin{enumerate}
\item \textbf{Q — Comprendre la question} : soulignez le verbe d'instruction (\emph{définir}, \emph{expliquer}, \emph{montrer}, \emph{discuter}) et l'objet de la question. Ne répondez pas à côté.
\item \textbf{R — Rédiger une réponse complète} : rédigez en phrases (jamais de réponse télégraphique), en commençant par reprendre les termes de la question.
\item \textbf{J — Justifier par le texte} : appuyez chaque réponse sur un passage précis du texte, cité entre guillemets et brièvement commenté. Une citation non expliquée ne vaut rien.
\end{enumerate}
\textbf{Réflexes} : une question = un paragraphe ; une citation courte (une ligne maximum) ; toujours annoncer la citation (« comme le dit l'auteur : ... »).
\end{methode}

\begin{exoresolu}[Répondre à une question guidée]
\textbf{Énoncé.} À partir du texte suivant : « La liberté n'est pas le pouvoir de faire ce que l'on veut, mais le pouvoir de vouloir ce que l'on doit. L'homme libre n'est pas celui qui obéit à ses caprices, car le caprice est un maître caché. » — Expliquez la distinction que l'auteur établit entre deux conceptions de la liberté.
\tcblower
\textbf{Solution détaillée.}\\
\textbf{Étape Q.} Le verbe d'instruction est « expliquez » : il faut dégager le sens de la distinction, pas simplement la citer. L'objet : deux conceptions opposées de la liberté.\\
\textbf{Étape R et J (rédaction).}\\
L'auteur oppose deux conceptions de la \cle{liberté}. La première, qu'il rejette, définit la liberté comme « le pouvoir de faire ce que l'on veut » : c'est la liberté conçue comme absence de contrainte, satisfaction immédiate de tous les désirs. La seconde, qu'il défend, définit la liberté comme « le pouvoir de vouloir ce que l'on doit » : c'est la liberté comme maîtrise de soi, accord de la volonté avec le devoir et la raison.\\
L'argument qui justifie ce choix est le suivant : celui qui obéit à ses caprices se croit libre, mais il est en réalité esclave, car « le caprice est un maître caché ». Autrement dit, l'obéissance aux désirs changeants est une servitude déguisée : le désir commande sans que nous en ayons conscience. La vraie liberté suppose donc de commander à ses désirs au lieu de leur obéir.\\
\textbf{Conclusion.} Pour l'auteur, la liberté n'est pas l'absence de loi, mais l'obéissance à la loi que l'on se donne rationnellement : être libre, c'est se gouverner soi-même.
\end{exoresolu}

\courssub{Rédiger et justifier une démarche, une réponse ou une conclusion}

\begin{methode}[Problématiser un texte : construire la question philosophique]
L'essai personnel exige de partir du texte pour poser un \emph{problème}. Démarche en 4 étapes :
\begin{enumerate}
\item \textbf{Relever la tension du texte} : tout texte philosophique oppose deux termes (liberté/nécessité, technique/homme, développement/tradition). Nommez cette opposition.
\item \textbf{Transformer la thèse en question} : si l'auteur affirme « P », demandez-vous « faut-il vraiment P ? » ou « P est-il toujours vrai ? ».
\item \textbf{Formuler la problématique} : une seule question, claire, philosophique (elle porte sur des concepts, pas sur des faits), qui appelle une discussion.
\item \textbf{Annoncer le plan} : l'essai personnel suit le plan en trois moments : thèse (ce que dit l'auteur), antithèse (objection, limite), synthèse (votre position justifiée).
\end{enumerate}
\textbf{Formule à retenir} : problématique = tension entre deux concepts + question ouverte.
\end{methode}

\begin{exoresolu}[Construire une problématique]
\textbf{Énoncé.} Texte : « Le développement d'un pays ne se mesure pas seulement à ses richesses, mais à la capacité de ses citoyens à penser par eux-mêmes. » — Dégagez la tension du texte et formulez une problématique pour un essai personnel.
\tcblower
\textbf{Solution détaillée.}\\
\textbf{Étape 1 — La tension.} Le texte oppose deux conceptions du \cle{développement} : le développement économique (les « richesses ») et le développement humain ou intellectuel (la « capacité à penser par soi-même »).\\
\textbf{Étape 2 — De la thèse à la question.} L'auteur affirme que le vrai développement est d'abord intellectuel et moral. Je peux alors demander : la richesse matérielle ne suffit-elle pas ? La pensée critique est-elle vraiment la condition du développement ?\\
\textbf{Étape 3 — Problématique.} « Le développement d'une nation repose-t-il d'abord sur ses richesses économiques ou sur la formation de l'esprit critique de ses citoyens ? »\\
\textbf{Étape 4 — Plan annoncé.}
\begin{enumerate}
\item \textbf{Thèse} : le développement est d'abord matériel (infrastructures, industries, revenus) — c'est le sens commun et la politique de nombreux États.
\item \textbf{Antithèse} : mais la richesse sans esprit critique est fragile — un peuple riche mais incapable de penser dépend des autres (exemple : importation de technologies sans maîtrise).
\item \textbf{Synthèse} : le développement véritable unit les deux dimensions, la richesse matérielle étant un moyen et l'autonomie de pensée la fin.
\end{enumerate}
\textbf{Conclusion.} La problématique est philosophique parce qu'elle interroge la \emph{nature} du développement, et non un fait particulier.
\end{exoresolu}

\begin{methode}[Rédiger l'introduction de l'essai personnel]
L'introduction se construit en 4 mouvements, dans l'ordre :
\begin{enumerate}
\item \textbf{L'amorce} : une phrase générale qui part du texte ou d'un constat de la vie courante (exemple camerounais bienvenu) et attire l'attention.
\item \textbf{L'analyse du sujet} : définissez les termes-clés du sujet en une ou deux phrases.
\item \textbf{La problématique} : posez la question philosophique construite à partir de la tension du texte.
\item \textbf{L'annonce du plan} : « Nous verrons d'abord que..., puis que..., enfin que... ».
\end{enumerate}
\textbf{Interdits} : pas de définition de dictionnaire recopiée, pas de « depuis la nuit des temps », pas de problématique absente ou noyée.
\end{methode}

\begin{exoresolu}[Rédiger une introduction complète]
\textbf{Énoncé.} Sujet : « La technique libère-t-elle l'homme ? » Rédigez l'introduction d'un essai personnel.
\tcblower
\textbf{Solution détaillée (rédaction intégrale).}\\
\textbf{Amorce} : Au Cameroun comme partout, le téléphone portable a transformé la vie quotidienne : on paie, on apprend, on travaille à distance. La technique semble avoir rendu l'homme plus puissant et plus libre.\\
\textbf{Analyse des termes} : La \cle{technique} désigne l'ensemble des outils et procédés inventés par l'homme pour transformer la nature ; la \cle{liberté} désigne le pouvoir d'agir selon sa volonté, sans contrainte ni servitude.\\
\textbf{Problématique} : Mais cette puissance nouvelle fait-elle de nous des êtres libres, ou crée-t-elle de nouvelles dépendances ? La technique libère-t-elle l'homme, ou risque-t-elle de l'asservir ?\\
\textbf{Annonce du plan} : Nous montrerons d'abord que la technique libère l'homme des contraintes naturelles ; nous verrons ensuite qu'elle engendre de nouvelles servitudes ; nous soutiendrons enfin que la technique ne libère que si elle reste un moyen maîtrisé par la raison.\\
\textbf{Conclusion de l'exercice} : cette introduction respecte les quatre mouvements obligatoires et annonce un plan dialectique (thèse, antithèse, synthèse), conforme aux exigences du Probatoire et du Baccalauréat technique.
\end{exoresolu}

\begin{methode}[Rédiger le développement : la règle A-E-L]
Chaque partie du développement suit la règle A-E-L :
\begin{enumerate}
\item \textbf{A — Affirmer} : commencez le paragraphe par une phrase-thèse claire (« La technique libère d'abord l'homme de la nature »).
\item \textbf{E — Expliquer et Exemplifier} : développez l'argument avec un raisonnement, puis illustrez-le par un exemple précis, de préférence tiré de la vie courante camerounaise (agriculture mécanisée à Mbanga, mobile money, forages d'eau).
\item \textbf{L — Lier} : terminez par une phrase de transition qui annonce la partie suivante (« Mais cette libération a un prix... »).
\end{enumerate}
\textbf{Réflexes} : un paragraphe = une idée ; citer brièvement un auteur étudié si possible (ex. : « comme le montre Bergson, l'intelligence fabrique des outils ») ; rédiger au présent de vérité.
\end{methode}

\begin{exoresolu}[Rédiger un paragraphe de développement]
\textbf{Énoncé.} Dans un essai sur le sujet « La technique libère-t-elle l'homme ? », rédigez le premier paragraphe de la première partie (thèse) en appliquant la règle A-E-L.
\tcblower
\textbf{Solution détaillée (rédaction intégrale).}\\
\textbf{A — Affirmation} : La technique libère d'abord l'homme des contraintes de la nature en augmentant sa puissance d'action.\\
\textbf{E — Explication et exemple} : En effet, l'homme est un être naturellement faible : sans outils, il ne peut ni labourer de grandes surfaces, ni se déplacer rapidement, ni conserver ses récoltes. La technique compense cette faiblesse. Ainsi, dans la région de Mbanga (Moungo), l'introduction du tracteur a permis aux agriculteurs de cultiver en quelques jours des surfaces qui exigeaient auparavant des semaines de travail manuel : le temps gagné est du temps libéré pour d'autres activités, comme l'éducation des enfants ou le commerce. De même, le forage hydraulique épargne aux femmes de nombreux villages des heures de marche quotidienne vers les points d'eau. La machine travaille à la place de l'homme et lui rend du temps et de l'énergie.\\
\textbf{L — Lien} : Cependant, cette libération n'est pas sans contrepartie : l'outil qui soulage peut aussi enchaîner, comme nous allons le voir.\\
\textbf{Conclusion de l'exercice} : le paragraphe respecte la structure A-E-L : une idée claire, un raisonnement, deux exemples camerounais précis, une transition vers l'antithèse.
\end{exoresolu}

\begin{methode}[Rédiger la conclusion de l'essai personnel]
La conclusion se rédige en 3 mouvements :
\begin{enumerate}
\item \textbf{Bilan} : résumez en deux phrases les positions confrontées (thèse et antithèse), sans copier les phrases du développement.
\item \textbf{Réponse à la problématique} : reprenez la question posée dans l'introduction et répondez-y clairement, en justifiant votre position personnelle.
\item \textbf{Ouverture} : terminez par une question nouvelle, liée au sujet, qui prolonge la réflexion (ex. : vers la responsabilité, l'éthique, l'éducation).
\end{enumerate}
\textbf{Interdits} : ne jamais introduire un argument nouveau ; ne jamais laisser la problématique sans réponse ; éviter « en conclusion » répété.
\end{methode}

\begin{exoresolu}[Rédiger une conclusion]
\textbf{Énoncé.} Rédigez la conclusion de l'essai sur le sujet « La technique libère-t-elle l'homme ? », sachant que le développement a montré : (1) la technique libère des contraintes naturelles ; (2) elle crée des dépendances nouvelles ; (3) elle ne libère que si elle est maîtrisée par la raison.
\tcblower
\textbf{Solution détaillée (rédaction intégrale).}\\
\textbf{Bilan} : Notre réflexion a d'abord montré que la technique accroît la puissance de l'homme et le délivre des servitudes naturelles ; elle a ensuite révélé que cette puissance se retourne parfois contre lui, lorsque l'outil devient un maître.\\
\textbf{Réponse à la problématique} : La technique libère-t-elle donc l'homme ? Nous répondons : oui, mais à une condition. En elle-même, la technique est un simple moyen, ni bon ni mauvais ; c'est l'usage que la raison en fait qui décide de sa valeur. L'homme n'est libéré par ses outils que s'il demeure capable de les juger et de s'en passer : la vraie libération est donc d'abord intérieure, elle suppose l'éducation de la volonté et de l'esprit critique.\\
\textbf{Ouverture} : Reste alors une question décisive pour nos sociétés : l'école forme-t-elle suffisamment des esprits capables de maîtriser les outils qu'elle met entre leurs mains ?\\
\textbf{Conclusion de l'exercice} : la conclusion répond explicitement à la problématique, prend position de façon nuancée et ouvre sur un problème voisin, conformément aux exigences du correcteur.
\end{exoresolu}

\begin{methode}[Justifier une réponse : citer et commenter le texte]
Pour qu'une réponse soit \emph{justifiée} et non seulement affirmée :
\begin{enumerate}
\item \textbf{Annoncer} la citation : « L'auteur affirme que... », « Comme le dit le texte... ».
\item \textbf{Citer} brièvement, entre guillemets, les mots exacts du texte (une ligne au maximum).
\item \textbf{Commenter} : expliquez en vos propres termes ce que la citation signifie et en quoi elle prouve votre réponse (« Cela signifie que... », « Autrement dit... »).
\item \textbf{Conclure} la réponse par une phrase qui revient explicitement sur la question posée.
\end{enumerate}
\textbf{Formule à retenir} : Citation + commentaire = preuve ; citation seule = paraphrase notée zéro.
\end{methode}

\begin{exoresolu}[Justifier une réponse par le texte]
\textbf{Énoncé.} Texte : « L'outil est bon serviteur, mais mauvais maître. Celui qui ne réfléchit pas à l'usage de ses instruments finit par leur obéir. » — Question : montrez, en vous appuyant sur le texte, que la technique exige une réflexion morale.
\tcblower
\textbf{Solution détaillée.}\\
\textbf{Annonce et citation} : L'auteur soutient que la technique n'est pas dangereuse par elle-même, mais par l'usage irréfléchi qu'on en fait. Il écrit : « L'outil est bon serviteur, mais mauvais maître. »\\
\textbf{Commentaire} : Cette formule signifie que l'outil est neutre : tant qu'il sert une fin choisie par l'homme, il est utile ; mais dès que l'homme cesse de décider et se laisse guider par l'outil, les rôles s'inversent et l'instrument commande. C'est ce que confirme la seconde phrase : « Celui qui ne réfléchit pas à l'usage de ses instruments finit par leur obéir. » Autrement dit, l'absence de réflexion transforme l'usage en dépendance : on pense par exemple à l'usager qui ne peut plus se passer de son téléphone et organise toute sa journée autour de lui.\\
\textbf{Conclusion} : Ainsi, le texte montre bien que la technique exige une réflexion morale : il faut sans cesse interroger la fin et les limites de nos outils pour qu'ils restent des serviteurs et ne deviennent pas des maîtres.
\end{exoresolu}

\courssub{TP Intégré : Simulation Bac — Corpus « Philosophie et Développement au Cameroun »}

\begin{exoresolu}[Sujet complet d'entraînement (3h) — Corrigé type]
\textbf{Énoncé.}\\
\textbf{Consignes générales} : L'épreuve comporte deux parties. La première partie (questions guidées) est notée sur 8 points. La seconde partie (essai personnel) est notée sur 12 points. La rédaction, la correction de la langue et la qualité du raisonnement entrent dans la notation.\\
\textbf{Corpus :}\\
« Le développement authentique ne saurait se réduire à la croissance quantitative des biens matériels. Il implique la promotion de l'homme tout entier et de tous les hommes. Or, dans nos pays, l'obsession du "rattrapage" technique nous fait souvent confondre les moyens et les fins. On importe des machines, on construit des routes, on érige des barrages, mais on néglige l'éducation de l'esprit critique, la formation du citoyen responsable, le respect des cultures locales. Une technique sans éthique, un développement sans liberté, ne font qu'aggraver la dépendance. Le vrai progrès est celui qui rend l'homme plus homme, c'est-à-dire plus libre et plus responsable. »\\
\textbf{Partie A : Questions guidées (8 points)}\\
1. Dégagez la thèse de l'auteur et la structure de son argumentation. (2 pts)\\
2. Expliquez la distinction que l'auteur opère entre « croissance » et « développement authentique ». (3 pts)\\
3. « Une technique sans éthique [...] ne fait qu'aggraver la dépendance. » Justifiez cette affirmation en vous appuyant sur le texte. (3 pts)\\
\textbf{Partie B : Essai personnel (12 points)}\\
Sujet : « Le développement d'une nation dépend-il d'abord de ses équipements techniques ou de la formation de ses citoyens ? »\\
\vspace{0.5cm}
\tcblower
\textbf{Corrigé détaillé.}\\
\textbf{Partie A : Questions guidées}\\
\textbf{1. Thèse et structure.}\\
\textbf{Thèse} : L'auteur soutient que le développement authentique n'est pas la simple accumulation de biens techniques (croissance), mais la promotion intégrale de l'homme (liberté, responsabilité, esprit critique) ; la technique sans éthique aggrave l'aliénation.\\
\textbf{Structure} : Trois moments. \textbf{Moment 1 (Définition)} : « Le développement authentique ne saurait se réduire... » — opposition croissance/développement. \textbf{Moment 2 (Constat critique)} : « Or, dans nos pays... » — confusion moyens/fins, importation technique sans formation humaine. \textbf{Moment 3 (Conclusion normative)} : « Le vrai progrès est celui... » — critère du progrès : l'humanisation (liberté, responsabilité).\\
\textbf{2. Distinction croissance / développement authentique.}\\
L'auteur oppose deux concepts. La \textbf{croissance} est quantitative : accumulation de biens, d'infrastructures, de machines (« importer des machines, construire des routes »). C'est une logique de « rattrapage » technique, purement matérielle. Le \textbf{développement authentique} est qualitatif : « promotion de l'homme tout entier », « éducation de l'esprit critique », « formation du citoyen responsable », « respect des cultures locales ». La croissance est un \emph{moyen} ; le développement est la \emph{fin} (l'homme plus libre, plus responsable).\\
\textbf{3. Justification : technique sans éthique = dépendance aggravée.}\\
\textbf{Annonce et citation} : L'auteur explique que l'obsession du rattrapage technique conduit à « confondre les moyens et les fins ». Il précise : « On importe des machines [...] mais on néglige l'éducation de l'esprit critique ».\\
\textbf{Commentaire} : Importer des machines sans former les esprits à les comprendre, les maîtriser et les adapter, c'est rester dépendant du fournisseur étranger. La technique devient alors une « boîte noire » que l'on subit. L'absence d'éthique (réflexion sur la fin, le bien commun, la justice) fait de la technique un outil de domination (dépendance économique, culturelle, politique) au lieu d'un outil d'émancipation. « Aggraver la dépendance » signifie : plus on a d'outils qu'on ne maîtrise pas, plus on est lié à ceux qui les produisent et les contrôlent.\\
\textbf{Conclusion partielle} : Le texte montre que l'autonomie technique suppose l'autonomie de la pensée.\\
\textbf{Partie B : Essai personnel — Éléments de corrigé}\\
\textbf{Problématique} : Le développement d'une nation dépend-il d'abord de ses équipements techniques (infrastructures, usines, numérique) ou de la formation de ses citoyens (éducation, esprit critique, éthique, culture) ?\\
\textbf{Plan dialectique :}\\
\textbf{I. Thèse : Les équipements techniques sont le socle matériel indispensable du développement.}\\
\begin{itemize}
\item Argument 1 : Satisfaction des besoins vitaux (eau, énergie, santé, transport). Ex. : Barrage de Lom Pangar (énergie), routes pour désenclaver les zones de production (Mbouda, Bamenda).
\item Argument 2 : Productivité et compétitivité économique. Usines de transformation locale (cacao, coton, bois) créent de la valeur ajoutée.
\item Argument 3 : Souveraineté technique. Maîtriser des infrastructures propres réduit la dépendance extérieure.
\end{itemize}
\textbf{II. Antithèse : Mais les équipements sans hommes formés sont stériles, voire dangereux (dépendance, gaspillage, aliénation).}\\
\begin{itemize}
\item Argument 1 : « White elephants » : infrastructures inutilisées ou mal entretenues par manque de compétences locales (ex. : certains équipements hospitaliers high-tech inopérants par absence de techniciens formés).
\item Argument 2 : Dépendance technologique : importer des machines sans transfert de compétences ni adaptation locale (ex. : agriculture mécanisée dépendante des pièces détachées et carburant importés).
\item Argument 3 : Risque éthique et social : technique sans régulation = corruption, destruction environnementale, exclusion sociale. Le citoyen non formé subit la technique au lieu de la piloter.
\end{itemize}
\textbf{III. Synthèse : Le développement véritable est la dialectique de l'outil et de l'esprit : l'équipement est le corps, la formation citoyenne est l'âme.}\\
\begin{itemize}
\item Argument 1 : Co-évolution nécessaire. Politique d'industrialisation \emph{adossée} à une réforme éducative massive (STEM, formation professionnelle, éthique de l'ingénieur).
\item Argument 2 : L'exemple des « Tigers asiatiques » (Corée du Sud) : investissement simultané dans l'école (capital humain) et l'usine (capital physique).
\item Argument 3 : Pour le Cameroun : la \cle{Stratégie Nationale de Développement 2020-2030 (SND30)} vise la transformation structurelle par le capital humain. Le développement durable (ODD 4, 8, 9) unit infrastructure résiliente et éducation de qualité.
\end{itemize}
\textbf{Introduction type (amorce + analyse + problématique + plan)} : Voir méthode ci-dessus (ex. téléphone portable / mobile money).\\
\textbf{Conclusion type (bilan + réponse + ouverture)} : « En définitive, les routes et les usines sont le squelette de la nation, mais l'école et la citoyenneté en sont l'esprit. Un pays bien équipé mais mal éduqué est un corps sans conscience. Le développement ne se décrète pas par décret d'importation, il se construit par l'éducation. \textbf{Ouverture} : Comment réformer concrètement nos systèmes éducatifs pour qu'ils produisent non seulement des techniciens, mais des citoyens ingénieurs de leur propre avenir ? »
\end{exoresolu}

\begin{attention}[Les 5 erreurs qui coûtent des points au Bac]
\begin{enumerate}
\item \textbf{Paraphraser le texte au lieu de l'expliquer.} Recopier le texte avec d'autres mots ne montre aucune compréhension. Exigez-vous de dégager la thèse et les arguments, et de les formuler avec \emph{vos} mots.
\item \textbf{Citer sans commenter.} Une citation jetée dans la copie sans explication est considérée comme du remplissage. Toute citation doit être annoncée, citée brièvement (entre guillemets) puis commentée.
\item \textbf{Oublier la problématique dans l'introduction.} Une introduction sans question philosophique claire condamne tout le devoir : le correcteur ne sait pas ce que vous discutez. La problématique doit être une question, formulée explicitement.
\item \textbf{Répondre à côté de la question posée.} Beaucoup d'élèves récitent un cours appris par cœur au lieu de répondre à la question précise. Soulignez le verbe d'instruction (\emph{définir}, \emph{expliquer}, \emph{discuter}) et vérifiez en fin de paragraphe que vous y avez répondu.
\item \textbf{Conclure sans répondre à la problématique.} Une conclusion qui se contente de résumer, ou qui introduit un argument nouveau, perd des points. La conclusion doit répondre clairement à la question de l'introduction, puis s'ouvrir sur une question nouvelle.
\end{enumerate}
\textbf{Réflexe final} : relisez chaque réponse en vous demandant : « Ai-je répondu à la question posée, et l'ai-je justifié par le texte ou par un argument ? »
\end{attention}

\newpage

\coursec{Exercices}

\c!!!!!!!!!!!!!!!!!!!!!!!!!!!!!!!!

\newpage

\coursec{Corrigés des exercices}

\begin{corrige}[Exercice 1 — QCM : repérer la nature de l'exercice sur texte]
\textbf{Question 1.} L'explication de texte philosophique vise principalement : \textbf{réponse b)} dégager la thèse de l'auteur, l'expliquer et la discuter. \emph{Justification :} il ne s'agit ni de paraphraser (réponse a), ni de donner d'emblée son opinion personnelle (réponse c) : la première tâche est de comprendre et de restituer fidèlement la pensée de l'auteur.

\textbf{Question 2.} La thèse d'un texte est : \textbf{réponse c)} l'idée principale que l'auteur défend et qu'il cherche à établir. On la formule en une phrase complète, à l'indicatif.

\textbf{Question 3.} Dans l'introduction d'une explication de texte, on trouve obligatoirement : \textbf{réponse a)} la présentation de l'auteur et du texte, la thèse, la problématique et l'annonce du plan. La biographie détaillée (réponse b) est un hors-sujet fréquent.

\textbf{Question 4.} L'essai personnel se distingue des réponses aux questions parce que : \textbf{réponse b)} il exige une prise de position argumentée et rédigée en continu par le candidat.

\textbf{Question 5.} La conclusion de l'essai personnel doit : \textbf{réponse c)} répondre clairement à la problématique et ouvrir éventuellement sur une question nouvelle.

\textbf{Points de vigilance :} au Probatoire technique, toute réponse de QCM non justifiée peut être pénalisée ; citez toujours la définition du cours qui fonde votre choix.
\end{corrige}

\begin{corrige}[Exercice 2 — Identifier thèse, arguments et structure d'un texte]
\textbf{1. La thèse.} Le texte proposé défend l'idée que le développement technique ne saurait se réduire à l'accumulation des richesses matérielles : il suppose un progrès de l'homme lui-même, de sa liberté et de sa rationalité. Formulation attendue en une phrase : \emph{« Le véritable développement est d'abord un développement humain et non seulement économique. »}

\textbf{2. Les arguments.} On relève trois arguments :
\begin{enumerate}
\item un argument par l'absurde : si le développement était seulement matériel, un pays riche serait automatiquement développé, ce que contredisent des situations concrètes (inégalités, dépendance) ;
\item un argument d'autorité implicite : le recours à l'idée que l'homme est « la fin et non le moyen » du développement ;
\item un argument par l'exemple : la référence à des pays qui, malgré leurs ressources (pétrole, minerais), demeurent en difficulté — situation que le candidat camerounais peut illustrer.
\end{enumerate}

\textbf{3. La structure.} Le texte procède en trois moments : constat (la confusion courante entre richesse et développement), critique (réfutation de cette confusion), ouverture (le développement comme épanouissement humain).

\textbf{Barème indicatif (sur 6 points) :} thèse correctement formulée : 2 pts ; deux arguments relevés et cités : 2 pts ; structure dégagée : 1 pt ; exactitude des citations : 1 pt.

\textbf{Points de vigilance :} ne pas confondre thèse et problématique ; citer le texte entre guillemets pour appuyer chaque réponse ; un argument sans référence textuelle n'est pas validé.
\end{corrige}

\begin{corrige}[Exercice 3 — Rédiger l'introduction d'une explication de texte]
\textbf{Modèle d'introduction attendu (en quatre temps) :}

\emph{Présentation :} Le texte soumis à notre étude est extrait d'une réflexion consacrée au rapport entre philosophie et développement. L'auteur y interroge la place de la pensée critique dans la transformation des sociétés africaines.

\emph{Thèse :} L'auteur y soutient que la philosophie, loin d'être un luxe inutile, est une condition du développement authentique, parce qu'elle forme des esprits capables de juger, d'innover et de résister aux préjugés.

\emph{Problématique :} Mais comment une discipline réputée abstraite peut-elle contribuer concrètement au développement d'un pays comme le Cameroun ?

\emph{Annonce du plan :} Nous expliquerons d'abord le sens et la portée de cette thèse ; nous en dégagerons ensuite les arguments ; nous la discuterons enfin pour en mesurer la validité.

\textbf{Barème indicatif (sur 4 points) :} présentation sobre de l'auteur et du texte : 0,5 pt ; thèse exacte en une phrase : 1 pt ; problématique sous forme interrogative : 1 pt ; annonce du plan en deux ou trois moments : 1 pt ; correction de la langue : 0,5 pt.

\textbf{Points de vigilance :} aucune biographie de l'auteur ; la problématique doit être une \emph{question}, pas une affirmation ; ne pas donner son avis dans l'introduction.
\end{corrige}

\begin{corrige}[Exercice 4 — Réponses aux questions guidées]
\textbf{Question 1 : Que signifie l'expression « esprit critique » dans le texte ?} Elle désigne l'aptitude à examiner les idées reçues, à demander des preuves avant d'adhérer, à distinguer le vrai du faux. Dans le texte, elle s'oppose à la soumission aux préjugés et à la rumeur. \emph{Attendu :} une définition en une ou deux phrases, appuyée sur un passage cité.

\textbf{Question 2 : Quelle est, selon l'auteur, la fonction de la philosophie dans la société ?} Sa fonction est formatrice et libératrice : elle forme le jugement des citoyens et les libère des opinions toutes faites, condition d'un développement maîtrisé. \emph{Attendu :} restitution fidèle, sans ajout personnel.

\textbf{Question 3 : L'auteur a-t-il raison d'opposer tradition et esprit critique ?} Réponse nuancée attendue : oui, si la tradition est vécue comme répétition aveugle (elle bloque alors l'innovation) ; non, si la tradition est transmise de manière réfléchie, car elle peut porter des valeurs (solidarité, parole donnée) utiles au vivre-ensemble. Le candidat doit prendre position et justifier par un exemple camerounais (palabre, conseil des notables, modernisation des coutumes).

\textbf{Barème indicatif (sur 6 points) :} 1,5 pt par question pour l'exactitude du contenu, 0,5 pt par question pour l'appui sur le texte ou l'exemple.

\textbf{Points de vigilance :} répondre à \emph{chaque} question dans l'ordre ; distinguer les questions de compréhension (rester fidèle à l'auteur) des questions de discussion (prendre position) ; toute affirmation sans justification perd la moitié des points.
\end{corrige}

\begin{corrige}[Exercice 5 — Essai personnel : « La philosophie peut-elle contribuer au développement du Cameroun ? »]
\textbf{Corrigé-type (plan détaillé rédigé).}

\emph{Introduction.} Amorce : on oppose souvent la philosophie, discipline des idées, au développement, affaire de techniques et d'économie. Problématique : la philosophie peut-elle contribuer au développement du Cameroun ? Annonce : nous verrons d'abord que la philosophie semble inutile au développement, puis qu'elle lui est en réalité indispensable, enfin qu'elle en est la condition critique et éthique.

\emph{Première partie (thèse de l'inutilité apparente).} La philosophie ne produit ni routes, ni usines, ni récoltes ; elle semble un luxe de pays riches. Argument : le développement exige des techniciens, des ingénieurs, des agronomes. Exemple : un chantier de construction à Douala n'attend pas de discours mais du béton.

\emph{Deuxième partie (antithèse : sa nécessité réelle).} Sans pensée critique, le développement est subi ou détourné : mauvais choix, corruption, imitation servile. La philosophie forme le jugement, apprend à problématiser, à décider en connaissance de cause. Exemples camerounais : l'émergence suppose des cadres capables d'analyser les politiques publiques ; la lutte contre les rumeurs et les préjugés (sur la santé, les vaccins) exige l'esprit critique.

\emph{Troisième partie (synthèse).} La philosophie ne remplace pas la technique : elle l'oriente. Elle est la condition éthique et critique d'un développement humain, qui mette l'homme au centre. Position personnelle argumentée : oui, la philosophie contribue au développement, non en produisant des biens, mais en formant des hommes libres et responsables.

\emph{Conclusion.} Réponse claire à la problématique ; ouverture : quelle place faire à l'enseignement de la philosophie dans les lycées techniques ?

\textbf{Barème indicatif (sur 8 points) :} introduction complète : 1,5 pt ; deux thèses opposées argumentées : 3 pts ; synthèse avec position personnelle justifiée : 1,5 pt ; exemples pertinents : 1 pt ; langue et présentation : 1 pt.

\textbf{Points de vigilance :} rédaction en continu (pas de titres apparents dans la copie) ; chaque partie contient au moins un argument et un exemple ; la position personnelle n'apparaît qu'en synthèse, jamais dès l'introduction.
\end{corrige}

\begin{corrige}[Exercice 6 — Simulation Bac : repérage et correction des erreurs fréquentes]
\textbf{1. Relever les erreurs de la copie présentée.} On attend l'identification d'au moins quatre fautes typiques :
\begin{enumerate}
\item la paraphrase : le candidat recopie le texte en changeant quelques mots au lieu de l'expliquer ;
\item le hors-sujet biographique : longue présentation de la vie de l'auteur sans lien avec la thèse ;
\item l'avis personnel prématuré : « je pense que » dès l'introduction ;
\item l'absence de problématique : l'introduction enchaîne directement sur le plan ;
\item les citations non signalées : passages du texte recopiés sans guillemets.
\end{enumerate}

\textbf{2. Correction de chaque erreur.}
\begin{enumerate}
\item Remplacer la paraphrase par l'explication : dégager les concepts clés, définir les termes, montrer l'enchaînement des idées ;
\item Réduire la présentation de l'auteur à une phrase utile (nom, œuvre, contexte du texte) ;
\item Réserver la position personnelle à la discussion ou à l'essai personnel ;
\item Construire une problématique : transformer la thèse en question (« Mais comment... ? », « Faut-il... ? ») ;
\item Mettre toute citation entre guillemets et la commenter aussitôt.
\end{enumerate}

\textbf{Barème indicatif (sur 6 points) :} 0,5 pt par erreur repérée (maximum 2 pts) ; 1 pt par correction pertinente (maximum 4 pts).

\textbf{Points de vigilance :} cet exercice entraîne à l'auto-évaluation ; au Probatoire, la relecture finale de la copie (cinq minutes) permet d'éliminer la plupart de ces fautes. Rappel : une copie sans problématique ne peut obtenir la moyenne en explication de texte.
\end{corrige}

\newpage

\coursec{Fiche bilan}

\begin{popfigure}
\begin{tikzpicture}[mindmap, grow cyclic, every node/.style={concept, font=\small},
root concept/.append style={concept color=popblue, font=\small\bfseries, minimum size=3.2cm},
level 1/.append style={level distance=3.6cm, sibling angle=60, minimum size=2.4cm, font=\footnotesize}]
\node[root concept]{L'exercice sur texte philosophique}
child[concept color=poporange]{ node {Généralités : nature et enjeux} }
child[concept color=popgreen]{ node {Étape A : explication de texte} }
child[concept color=poppurple]{ node {Étape B : réponses guidées} }
child[concept color=poppink]{ node {Étape C : essai personnel} }
child[concept color=popgold]{ node {TP : simulation Bac} }
child[concept color=popblueL]{ node {Erreurs et remédiation} };
\end{tikzpicture}
\legende{Carte mentale de la leçon : les six étapes-clés pour réussir l'exercice sur texte philosophique au Probatoire et au Baccalauréat technique.}
\end{popfigure}

\begin{bilan}
\textbf{Définitions clés}
\begin{itemize}
\item \cle{Explication de texte} : exercice qui consiste à dégager le sens d'un texte philosophique, à en expliquer la thèse et les arguments, sans le paraphraser.
\item \cle{Thèse} : position défendue par l'auteur, réponse qu'il apporte à un problème.
\item \cle{Problématique} : question centrale à laquelle le texte répond ; elle guide toute l'étude.
\item \cle{Argument} : raison avancée par l'auteur pour justifier sa thèse.
\item \cle{Essai personnel} : réflexion argumentée du candidat qui discute la thèse de l'auteur (accord, désaccord, dépassement) à partir d'exemples concrets.
\end{itemize}

\textbf{À connaître par c\oe ur}
\begin{center}
\begin{tabularx}{\linewidth}{|l|X|}
\hline
\rowcolor{popblueL} \textbf{Élément} & \textbf{Contenu attendu} \\
\hline
Introduction & Présenter l'auteur et le texte, dégager la problématique, annoncer le plan. \\
\hline
Analyse & Expliquer la thèse, dégager les arguments, définir les notions-clés. \\
\hline
Discussion & Confronter la thèse à d'autres points de vue, prendre position. \\
\hline
Conclusion & Bilan de l'étude, réponse claire à la problématique, ouverture éventuelle. \\
\hline
\end{tabularx}
\end{center}

\textbf{Méthodes express}
\begin{itemize}
\item \emph{Lecture} : lire le texte 2 fois, souligner la thèse et les articulations logiques (donc, car, mais, cependant).
\item \emph{Questions guidées} : répondre dans l'ordre, en citant le texte, sans recopier de longs passages.
\item \emph{Essai personnel} : problématiser $\rightarrow$ confronter $\rightarrow$ argumenter avec des exemples camerounais concrets $\rightarrow$ conclure.
\item \emph{Gestion du temps} : environ 1/3 pour l'explication, 1/3 pour les questions, 1/3 pour l'essai.
\end{itemize}
\end{bilan}

\begin{aretenir}[Checklist avant le Bac]
\begin{itemize}
\item[$\square$] Je sais identifier la thèse de l'auteur en une phrase.
\item[$\square$] Je sais formuler la problématique sous forme de question.
\item[$\square$] Je distingue clairement thèse, argument et exemple.
\item[$\square$] Je définis les notions-clés du texte avant de discuter.
\item[$\square$] Je réponds aux questions dans l'ordre et en m'appuyant sur le texte.
\item[$\square$] Je cite brièvement le texte sans le paraphraser.
\item[$\square$] Je sais construire un essai personnel en trois temps (thèse, antithèse, synthèse).
\item[$\square$] J'illustre mes arguments avec des exemples concrets de la vie courante au Cameroun.
\item[$\square$] Je rédige une introduction avec problématique et annonce du plan.
\item[$\square$] Je conclus en répondant clairement à la problématique posée.
\item[$\square$] Je respecte le temps imparti et je relis ma copie.
\item[$\square$] Je soigne l'orthographe, la présentation et les articulations logiques.
\end{itemize}
\end{aretenir}

\findecours

\end{document}
